% ============================================================================
%   Claim safe AAAI Draft for the DST-EDL Architecture
%   Uncertainty to Topology: Adaptation for Structured Sparse Evidential Learning under Extreme Class Imbalance
%
%   File: dst_edl_report_en.tex
%   Draft label: tradeoff aware reported configuration (06/2026)
%   Language: English
% ============================================================================

\documentclass[letterpaper]{article} % DO NOT CHANGE THIS
\usepackage{aaai}  % DO NOT CHANGE THIS
\usepackage{times}  % DO NOT CHANGE THIS
\usepackage{helvet}  % DO NOT CHANGE THIS
\usepackage{courier}  % DO NOT CHANGE THIS
\usepackage[hyphens]{url}  % DO NOT CHANGE THIS
\usepackage{graphicx} % DO NOT CHANGE THIS
\urlstyle{rm} % DO NOT CHANGE THIS
\def\UrlFont{\rm}  % DO NOT CHANGE THIS
\frenchspacing  % DO NOT CHANGE THIS
\setlength{\pdfpagewidth}{8.5in} % DO NOT CHANGE THIS
\setlength{\pdfpageheight}{11in}
\setcounter{secnumdepth}{2}
%%%%%%%%%%
% PDFINFO for PDFLATEX
\pdfinfo{
/Title (Uncertainty to Topology: Adaptation for Structured Sparse Evidential Learning under Extreme Class Imbalance)
/Author (Anonymous Authors)
}
%%%%%%%%%%
 % DO NOT CHANGE THIS

% --- Additional Packages ---
\usepackage[utf8]{inputenc}
\usepackage[english]{babel}
\usepackage{amsmath, amssymb, amsthm, bm}
\usepackage{booktabs}
\usepackage{multirow}
\usepackage{array}
\usepackage{adjustbox}
\usepackage{placeins}
\usepackage{stfloats}
\usepackage{float}
\usepackage{enumitem}
\usepackage[final]{microtype}
\usepackage{tikz}
\usetikzlibrary{arrows.meta, positioning, calc}
\usepackage[table]{xcolor}
\usepackage{colortbl}
\usepackage[colorlinks=true,linkcolor=black,citecolor=blue,urlcolor=blue]{hyperref}

\newcolumntype{C}{>{\centering\arraybackslash}c}

% --- Custom commands ---
\newcommand{\Dir}{\text{Dir}}
\newcommand{\KL}{\text{KL}}
\newcommand{\R}{\mathbb{R}}
\newcommand{\E}{\mathbb{E}}
\newcommand{\Norm}{\text{Norm}}
\newcommand{\softplus}{\text{Softplus}}
\newcommand{\pospart}[1]{\left[#1\right]_+}
\DeclareMathOperator*{\argmax}{arg\,max}
\DeclareMathOperator*{\argmin}{arg\,min}
\DeclareMathOperator{\TopTwo}{Top2}

\tolerance=1500
\hyphenpenalty=50
\exhyphenpenalty=50
\lefthyphenmin=2
\righthyphenmin=3
\emergencystretch=1em
\hyphenation{ev-i-den-tial Di-rich-let un-cer-tain-ty clas-si-fi-ca-tion cal-i-bra-tion struc-tur-al top-ol-o-gy spec-i-fic-i-ty sen-si-tiv-i-ty crys-tal-li-za-tion rep-re-sen-ta-tion-al im-bal-ance com-pat-i-bil-i-ty ac-cel-er-a-tion de-ploy-ment fea-si-bil-i-ty mor-phol-o-gy con-cen-tra-tion reg-u-lar-i-za-tion}

\hbadness=10000
\vbadness=10000
\hfuzz=10pt
\vfuzz=4pt
\raggedbottom
\setlength{\textfloatsep}{10pt plus 2pt minus 3pt}
\setlength{\dbltextfloatsep}{10pt plus 2pt minus 3pt}
\setlength{\floatsep}{8pt plus 2pt minus 2pt}
\setlength{\dblfloatsep}{8pt plus 2pt minus 2pt}

\newtheorem{remark}{Remark}

% Compatibility patch for manual \bibitem entries with aaai.sty
\makeatletter
\renewenvironment{thebibliography}[1]{%
  \section*{References\@mkboth{REFERENCES}{REFERENCES}}%
  \list{[\arabic{enumiv}]}{%
    \settowidth\labelwidth{[#1]}%
    \leftmargin\labelwidth
    \advance\leftmargin\labelsep
    \usecounter{enumiv}%
    \let\p@enumiv\@empty
    \renewcommand\theenumiv{\arabic{enumiv}}%
    \itemsep .01in
  }%
  \sloppy\clubpenalty4000\widowpenalty4000\sfcode`\.\@m
}{\def\@noitemerr{\@latex@warning{Empty `thebibliography' environment}}\endlist}
\makeatother


% ============================================================================
\title{Uncertainty to Topology: Adaptation for Structured Sparse Evidential Learning under Extreme Class Imbalance}
\author{Anonymous Authors}

\begin{document}

\maketitle

\begin{abstract}
Extremely imbalanced classification requires rare positive ranking,
uncertainty estimation, and compatibility with structured sparsity.
Evidential Deep Learning (EDL) provides Dirichlet vacuity and ambiguity
signals in one forward pass. Dynamic Sparse Training (DST) adapts
network connectivity during optimization. Existing sparse update rules
usually use magnitude or task gradients as structural saliency. They do
not directly use the evidential state. We propose DST-EDL, a sparse
evidential framework that maps Dirichlet uncertainty to topology
selection pressure under exact valid block 2:4 sparsity. The method uses
a two time scale training rule. The prediction loop optimizes masked
weights with an evidential focal loss. The structural loop prunes active
links with a signed first order ambiguity over vacuity risk ratio. It
regrows dormant links with KL to uniform evidential saliency from a
virtual effective weight probe. Blockwise top two projection restores the
2:4 feasible set after each refresh. We verify positive Dirichlet states,
bounded uncertainty signals, masked gradient isolation, dormant link
saliency, bounded latent scores, and exact sparse projection. On ISIC
2024 (${\sim}0.15\%$ malignant prevalence), DST-EDL achieves the highest
mean pAUC$_{0.80}{=}0.1168$ and Macro AUROC${=}0.8891$ among the compared
dense evidential and sparse 2:4 baselines over seeds 42, 123, and 456.
The same runs preserve 2,789,376/2,789,376 valid sparse blocks. No
structural collapse is detected. The supported claim is specific:
DST-EDL improves high sensitivity ranking under exact structured
sparsity. It does not dominate calibration, PR AUC, or every
classification metric.
\end{abstract}

\section{Introduction}\label{sec:intro}

Extremely imbalanced classification arises in medical screening, industrial inspection, and fraud detection. False negatives on rare positives can be costly. Models in these settings require rare event sensitivity. They also require calibrated uncertainty estimates and deployment compatible inference. Structured sparsity adds a separate design constraint. The sparse model must decide which active connections remain useful for rare and uncertain regions.

Long tailed learning methods modify losses, logits, sampling rules, or classifier heads. They do not use uncertainty to adapt sparse topology. Evidential Deep Learning (EDL)~\cite{sensoy2018evidential} parameterizes a Dirichlet distribution over class probabilities. It exposes vacuity and expected entropy in one forward pass~\cite{malinin2018predictive}. Standard EDL models usually keep a dense or fixed topology predictor. Dynamic Sparse Training (DST)~\cite{evci2020rigging} adapts connectivity during optimization. Its structural criteria usually rely on magnitude or task gradients. NVIDIA 2:4 structured sparsity supplies a deployment compatible constraint. It does not define uncertainty aware mask adaptation.

We introduce \textbf{DST-EDL}, an uncertainty guided sparse evidential learning framework. The method treats evidential uncertainty as a topology selection signal under a fixed sparse budget. The design is summarized by the following constrained decomposition:
\begin{equation}
\begin{aligned}
\mathcal{D}_{\mathrm{sig}}(W,M)
&=
\mathbb{E}_{(x,y)}
\!\left[
L_{\mathrm{EFL}}(W\odot M;x,y)
\right] \\
&\quad
+ \lambda_s\Phi_{\mathrm{struct}}(W,M),
\qquad
M\in\mathcal{C}_{2:4}.
\end{aligned}
\end{equation}
We use \(\mathcal{D}_{\mathrm{sig}}\) only as a conceptual decomposition. It separates the ordinary training loss from the structural proxy quantities. It is not a Lyapunov function. It is not a descent certificate or a block coordinate objective. It is not minimized directly end to end. The contribution is the coupling between evidential uncertainty and topology evolution. The structural container is
\begin{equation}
\Phi_{\mathrm{struct}}(W,M) = \Phi_{\mathrm{prune}}(W,M) + \lambda_g \Phi_{\mathrm{grow}}(W,M),
\end{equation}
where the pruning and regrowth proxies are:
\begin{equation}
\Phi_{\mathrm{prune}} = \mathbb{E} \left[ \frac{u_a}{u_e+\epsilon} \right], \quad \Phi_{\mathrm{grow}} = \mathbb{E} [ \mathrm{KL}(\mathrm{Dir}(\alpha)\Vert \mathrm{Dir}(\mathbf{1})) ].
\end{equation}
The two terms are not optimized as a single penalty. DST-EDL uses them to define separate local saliency scores. The pruning score ranks active links by the predicted effect of removal on an ambiguity over vacuity proxy. The regrowth score ranks dormant links by their local influence on a KL uniform evidence state. The refreshed scores are projected exactly onto the valid block 2:4 feasible set. The reported model uses topology caching, softened signed first order pruning, KL uniform regrowth, symmetric KL regularization, and evidential focal learning. ISIC 2024 is the primary rare event case study. Broader validation is reported only when complete evidence is available.

\paragraph{Main Contributions.}
\begin{itemize}
    \item \textbf{Uncertainty guided structured topology adaptation.} We convert Dirichlet vacuity and expected entropy into structural pressure for pruning and regrowth under exact valid block 2:4 sparsity over eligible complete 4-blocks.
    \item \textbf{Decoupled sparse evidential training.} We separate ordinary masked prediction updates from amortized structural refresh, enabling dormant links to receive proxy saliency without ordinary masked parameter updates.
    \item \textbf{Conservative rare event evaluation.} We study the method on ISIC 2024 with repeated seeds, dense and sparse baselines, ranking metrics, calibration metrics, and structural validity diagnostics. The claims are restricted to the reported evidence.
\end{itemize}

\section{Related Work}\label{sec:related}

\paragraph{Long Tailed and Extreme Imbalanced Learning.} Long tailed recognition is commonly addressed through resampling, reweighting, margin adjustment, classifier decoupling, calibration aware smoothing, and contrastive learning. Representative methods include Class Balanced Loss~\cite{cui2019class}, LDAM-DRW~\cite{cao2019learning}, Logit Adjustment~\cite{menon2021long}, Balanced Softmax~\cite{ren2020balanced}, decoupled classifier training~\cite{kang2019decoupling}, MiSLAS~\cite{zhong2021improving}, BBN~\cite{zhou2020bbn}, and PaCo~\cite{cui2021parametric}. These methods modify the loss, logits, sampling distribution, or classifier structure. DST-EDL studies a different mechanism. It uses evidential structural signals to adapt sparse topology under extreme imbalance.

\paragraph{Evidential and Dirichlet Based Uncertainty Estimation.} Evidential Deep Learning~\cite{sensoy2018evidential} replaces point estimate softmax confidence with a Dirichlet distribution over class probabilities. Related work includes Prior Networks~\cite{malinin2018predictive}, Posterior Networks~\cite{charpentier2020posterior}, Fisher EDL~\cite{deng2023fisher}, R-EDL~\cite{chen2023redl}, Re-EDL~\cite{chen2024reedl}, Flexible EDL~\cite{yoon2025flexible}, ANEDL~\cite{yu2024anedl}, and Evidence Contraction~\cite{wu2024evidence}. Recent critiques~\cite{shen2024mirage} caution against interpreting EDL uncertainty as exact Bayesian posterior uncertainty. Accordingly, DST-EDL uses Dirichlet vacuity and expected entropy as structural adaptation signals. These signals are not Bayesian uncertainty guarantees.

\paragraph{Calibration, Selective Prediction, and Failure Detection.} Rare event deployment requires calibrated confidence, operating point awareness, and selective abstention. Temperature scaling~\cite{guo2017calibration}, selective classification~\cite{geifman2017selective,geifman2019selectivenet}, AURC~\cite{ding2020revisiting}, and recent evaluation critiques~\cite{jaeger2023call,traub2024overcoming} motivate joint reporting of calibration and working point metrics. DST-EDL reports thresholded operating points, overall ECE, ranking metrics, and risk coverage behavior.

\paragraph{Dynamic Sparse Training and Structured Sparsity.} Dynamic sparse training maintains sparse connectivity during optimization~\cite{han2016deep,frankle2019lottery}. SNIP~\cite{lee2019snip}, SET~\cite{mocanu2018scalable}, SNFS~\cite{dettmers2019sparse}, RigL~\cite{evci2020rigging}, PFFDST~\cite{li2025pffdst}, and SRigL~\cite{lasby2023srigl} motivate sensitivity, regrowth, and structured sparsity. NVIDIA Sparse Tensor Cores motivate 2:4 deployment constraints~\cite{mishra2021accelerating}. Recent 2:4 systems also show that measured acceleration depends on tensor shapes, layout conversion, and backend kernel support rather than mask validity alone~\cite{hu2024accelerating,nvidiacusparselt,pytorchsemistructured}. DST-EDL refreshes the 2:4 mask from Dirichlet derived pruning and regrowth saliency.

\paragraph{Positioning of DST-EDL.} Prior work addresses complementary components. EDL models Dirichlet predictive uncertainty. DST adapts sparse connectivity through magnitude or task gradient saliency. Structured sparsity enforces hardware compatible sparse patterns. SNIP~\cite{lee2019snip} provides one shot connection sensitivity. SET~\cite{mocanu2018scalable} rewires sparse connectivity. SNFS~\cite{dettmers2019sparse} and RigL~\cite{evci2020rigging} use training gradient information for dynamic sparse regrowth. SRigL~\cite{lasby2023srigl} extends this line toward structured sparsity. Fisher EDL~\cite{deng2023fisher}, R-EDL~\cite{chen2023redl}, and Flexible EDL~\cite{yoon2025flexible} refine evidential uncertainty modeling. DST-EDL uses Dirichlet derived signals as direct structural pressure under exact valid block 2:4 feasibility.

\section{Preliminaries: Dirichlet Evidential Signals}\label{sec:math}

For logits $\bm{z}\in\R^K$, DST-EDL maps each class logit to non negative evidence using Softplus,
\begin{equation}\label{eq:softplus}
    e_c=\ln(1+\exp z_c),\qquad
    \alpha_c=e_c+1,\qquad
    S=\sum_{c=1}^K\alpha_c .
\end{equation}
Softplus is used instead of an output ReLU. ReLU can create zero gradient regions for negative evidence logits. The predictive mean is $\hat p_c=\alpha_c/S$. With uniform Subjective Logic base rate $a_c=1/K$, this equals
\begin{equation}\label{eq:p_hat}
    \hat p_c=b_c+a_cu_e=\frac{e_c+1}{S},
    \qquad b_c=\frac{e_c}{S},\qquad u_e=\frac{K}{S}.
\end{equation}
Imbalance is handled through the loss. The Dirichlet prior is unchanged. The identity $\alpha_c=e_c+1$ and the lower bound $S\ge K$ are preserved.

Following recent critiques of Evidential Deep Learning~\cite{shen2024mirage}, we do not interpret EDL signals as exact Bayesian posterior uncertainty. We use Dirichlet vacuity $u_e=K/S$ as a low evidence structural proxy and expected categorical entropy
\begin{equation}\label{eq:u_a}
    u_a = \sum_{c=1}^K \frac{\alpha_c}{S} \left[\psi(S + 1) - \psi(\alpha_c + 1)\right]
\end{equation}
as an ambiguity proxy. In implementation, $u_a$ is clamped below at zero to guard against floating point error.

For finite logits, Softplus gives $S>K$. The limiting lower bound $S \to K$ is approached as evidence tends to zero. This parameterization guarantees positive Dirichlet parameters, bounded vacuity, and a well defined predictive mean. Formal validity statements and proofs are provided in Appendix~\ref{app:proofs}.

\paragraph{Theoretical Scope.}
The quantities $u_e$ and $u_a$ are well defined Dirichlet derived signals. The adopted parameterization gives bounded $u_e$ and non negative $u_a$. These signals drive topology adaptation and uncertainty aware evaluation. Appendix~\ref{app:proofs} gives the proofs.


\section{DST-EDL Framework}\label{sec:methodology}

For exposition, we refer to the pruning signal as \emph{risk ratio signed pruning}, the regrowth signal as \emph{KL uniform structural regrowth}, and the cached refresh schedule as a \emph{two time scale topology cache}.

DST-EDL separates training into two coupled loops. The prediction loop updates weights at mini batch frequency. The structural loop updates topology at refresh frequency. This schedule reduces topology noise. It still permits dormant links to receive proxy structural saliency. The prediction loop uses a cached 2:4 mask $\bm{M}$. It optimizes active weights $\bm{W}$ with the evidential objective. The structural loop uses an amortized proxy batch. It computes $u_e,u_a$, forms $R_{\mathcal B}$ and $L_{\text{grow}}$, updates latent scores $\bm{V}$, and refreshes $\bm{M}$ by blockwise top two projection.

The reported model employs topology caching, softened signed first order pruning with pruning strength $0.5$, pruning warmup, KL uniform regrowth, symmetric KL regularization, and evidential focal learning. This configuration is conservative. It is not tuned as an aggressive sensitivity maximizer.

\subsection{2:4 Mask Constrained Sparse Topology}
DST-EDL wraps convolutional kernels and linear projections with binary masks. Prediction weights $\bm{W}^{(l)}$ are optimized by the ordinary training loss. Latent scores $\bm{V}^{(l)}$ determine the mask $\bm{M}^{(l)}$:
\begin{equation}\label{eq:wrapped_weights}
    \bm{W}_{\text{eff}}^{(l)}=\bm{W}^{(l)}\odot\bm{M}^{(l)} .
\end{equation}
For each eligible row and complete contiguous 4-block, exactly two entries are active:
\begin{equation}\label{eq:c24}
    \mathcal{C}_{2:4}
    =
    \left\{
    \bm{M}\in\{0,1\}^{d_{\text{out}}\times d_{\text{in}}}:
    \sum_{j=0}^{3} M_{i,4k+j}=2,\ \forall i,k
    \right\}.
\end{equation}
Convolutional tensors are flattened row wise before block partitioning. Incomplete trailing segments are excluded from the exact 2:4 claim. Given latent scores, projection is independent within each valid block:
\begin{equation}\label{eq:mask}
    \begin{aligned}
    I_{i,k}^{(l)} &= \TopTwo_{j\in\{0,1,2,3\}}
    V_{\text{flat},i,4k+j}^{(l)},\\
    M_{\text{flat},i,4k+j}^{(l)} &=
    \begin{cases}
        1 & \text{if } j\in I_{i,k}^{(l)},\\
        0 & \text{otherwise.}
    \end{cases}
    \end{aligned}
\end{equation}
This is equivalent to $\bm{M}^{(t+1)}=\Pi_{\mathcal{C}_{2:4}}(\bm{V}^{(t)})$. The projection maximizes $\langle \bm{M},\bm{V}\rangle$ over each complete 4-block. Every refreshed eligible block remains exactly 2:4 feasible by construction. Appendix~\ref{app:24_feasibility} gives the proof.

\subsection{Uncertainty Guided Pruning and KL uniform Structural Regrowth}
At each structural refresh, DST-EDL evaluates a proxy batch and computes two saliency signals. The pruning signal uses the ambiguity over vacuity ratio. The regrowth signal uses KL to uniform evidential influence. Ordinary prediction uses $\bm{W}\odot\bm{M}$. Structural scoring uses an auxiliary dense effective weight probe $\widetilde{\bm{W}}_{\text{eff}}$ initialized at the current masked values:
\begin{equation}\label{eq:virtual_probe}
\begin{aligned}
    \widetilde{\bm{W}}_{\text{eff}}
    &\leftarrow \bm{W}\odot\bm{M},\\
    \nabla_{\widetilde{\bm{W}}_{\text{eff}}}(\cdot)
    &\text{ is used only for mask scoring.}
\end{aligned}
\end{equation}
This probe gives dormant coordinates structural saliency without applying ordinary parameter updates to dormant weights. The formal distinction from masked training gradients is given in Appendix~\ref{app:virtual_probe_note}.

For pruning, the proxy batch risk ratio is
\begin{equation}\label{eq:batch_risk}
    R_{\mathcal{B}}
    =
    \frac{1}{|\mathcal{B}|}
    \sum_{n\in\mathcal{B}}
    \frac{u_a^{(n)}}{u_e^{(n)}+\epsilon}.
\end{equation}
The ratio is a design proxy with three desiderata. It increases with ambiguity. It attenuates pruning pressure under vacuity. It is bounded on finite proxy batches. We do not claim theoretical superiority over entropy, margin, Fisher style, or task gradient signals. For active links, signed first order pruning uses
\begin{equation}\label{eq:C_ij}
    C_{ij}
    =
    \operatorname{Rank}_{l}\!\left(
    \pospart{w_{\text{eff},ij}
    \frac{\partial R_{\mathcal B}}{\partial w_{\text{eff},ij}}}
    \right),
    \qquad
    \pospart{x}=\max(x,0).
\end{equation}
Here $\operatorname{Rank}_{l}$ is a layer wise percentile rank. The positive part keeps only removals locally predicted to reduce the proxy ratio. Appendix~\ref{app:signed_pruning} gives the Taylor interpretation and rank scale invariance.

For regrowth, the reported setting uses KL uniform regrowth:
\begin{equation}\label{eq:grow_objective}
    L_{\text{grow}}
    =
    \frac{1}{|\mathcal{B}|}\sum_{n\in\mathcal{B}}
    \KL\!\left(
    \Dir(\bm{\alpha}^{(n)})\middle\|\Dir(\mathbf{1})
    \right).
\end{equation}
Its absolute probe gradient rank defines dormant link influence:
\begin{equation}\label{eq:G_ij}
    G_{ij}
    =
    \operatorname{Rank}_{l}\!\left(
    \left|
    \frac{\partial L_{\text{grow}}}
    {\partial w_{\text{eff},ij}}
    \right|
    \right).
\end{equation}
This is an influence score. It is not a descent direction. It does not instruct the model to increase confidence. It asks which dormant coordinates could most affect the evidential state relative to the uninformative Dirichlet reference. The supervised loss, top two projection, and calibration protocol determine the final effect.

The two rank normalized scores are combined in the latent score update
\begin{equation}\label{eq:score_update}
\begin{aligned}
    A_{ij}^{(t)} &=
    G_{ij}^{(t)}-\lambda_p C_{ij}^{(t)},\\
    D_{ij}^{(t)} &=
    \rho G_{ij}^{(t)}+\xi_{ij}^{(t)},\\
    \Delta V_{ij}^{(t)}
    &=
    M_{ij}^{(t)}A_{ij}^{(t)}
    +\left(1-M_{ij}^{(t)}\right)D_{ij}^{(t)},\\
    \mu_{ij}^{(t)}
    &= \beta_m\mu_{ij}^{(t - 1)}
    +(1-\beta_m)\Delta V_{ij}^{(t)},\\
    X_{ij}^{(t)} &= V_{ij}^{(t - 1)}+\eta_{\text{struct}}\mu_{ij}^{(t)},\\
    Y_{ij}^{(t)} &= X_{ij}^{(t)}-\operatorname{mean}(\bm{X}^{(t)}),\\
    V_{ij}^{(t)} &= \operatorname{clamp}\!\left(
    Y_{ij}^{(t)},-V_{\max},V_{\max}\right).
\end{aligned}
\end{equation}
For active links, $G_{ij}^{(t)}$ acts as a retention and influence score.
$C_{ij}^{(t)}$ supplies signed removal pressure. For dormant links, the
attenuated regrowth score $\rho G_{ij}^{(t)}$ and optional tie breaking noise
$\xi_{ij}^{(t)}$ control reactivation. The final clamp bounds the latent
morphology state, and Eq.~\ref{eq:mask} restores exact valid block 2:4
feasibility after each refresh.

\subsection{Gradient Isolation, Caching, and Stability Diagnostics}
Ordinary training uses a fixed cached mask for each mini batch. Prediction gradients and weight decay update only active links. For any differentiable mini batch loss evaluated through $\bm{W}_{\text{eff}}=\bm{W}\odot\bm{M}$,
\begin{equation}\label{eq:masked_gradient}
    \nabla_{\bm{W}} L
    =
    \nabla_{\bm{W}_{\text{eff}}} L \odot \bm{M}.
\end{equation}
Dormant parameters receive zero ordinary task gradient under a fixed mask. DST-EDL separates this training path from the virtual structural probe in Eq.~\ref{eq:virtual_probe}. The probe supplies saliency for future regrowth. It is not applied as a dormant weight update.

\paragraph{Summary of Well Posedness and Algorithmic Validity.}
This section establishes algorithmic validity. It verifies valid Dirichlet derived signals, exact 2:4 feasibility over eligible complete blocks, local first order pruning and regrowth interpretations, bounded latent scores, virtual probe dormant saliency, and ordinary masked gradient isolation. These checks make the controller well defined and mask feasible. They do not establish convergence, optimality, Bayesian correctness, validation improvement, or hardware speedup. Appendix~\ref{app:proofs} and Table~\ref{tab:guarantees} state the boundary explicitly.

\section{Optimization Framework}\label{sec:optimization}

To address extreme class imbalance, we build on Evidential Focal Loss (EFL) from Zhou et al.~\cite{zhou2023domain}, which combines evidential Dirichlet classification with focal modulation. The reported DST-EDL configuration uses square root class frequency dampening, stop gradient focal modulation, KL annealing, and symmetric KL regularization.

In the reported symmetric setting, the KL multiplier is exactly uniform across samples and classes. The asymmetric KL multiplier is retained only for the diagnostic ablation named ``asymmetric KL''. It is not used in the reported DST-EDL row. The implemented reported objective is
\begin{equation}\label{eq:efl}
\begin{aligned}
    L_{\text{EFL}}
    &= \frac{1}{N}\sum_{n=1}^N \mathcal{F}_n(t)
    \Bigl[
    \bar{\omega}_{y_n}L_{\text{CE}}^{(n)} \\
    &\hspace{2.4em}
    + \lambda_{\text{KL}}\mu(t)
    L_{\text{KL}}^{(n)}
    \Bigr].
\end{aligned}
\end{equation}
The dampened class weights are normalized so that the majority class has weight one and then capped:
\begin{equation}\label{eq:sample_weight}
\begin{aligned}
    \omega_c
    &= \sqrt{\frac{N_{\text{majority}}}{N_c}},
    &\bar{\omega}_c
    &= \min(\omega_c,\omega_{\max}),\\
    \omega_{\text{majority}}
    &= 1,
    &\omega_{\max}
    &= 10.0 .
\end{aligned}
\end{equation}
The square root form damps inverse frequency scaling relative to direct inverse prevalence weighting. The cap bounds the imbalance modulation. The expected cross entropy under the Dirichlet predictive distribution is:
\begin{equation}\label{eq:dirichlet_ce}
    L_{\text{CE}}^{(n)}
    = \sum_{c=1}^K y_c^{(n)}
    \left[\psi(S^{(n)})-\psi(\alpha_c^{(n)})\right].
\end{equation}
The focal term is applied with stop gradient confidence so that focal weighting changes sample emphasis without directly distorting the Dirichlet evidence:
\begin{equation}\label{eq:focal_schedule}
\begin{aligned}
    \mathcal{F}_n(t)
    &= \left(1-\operatorname{sg}\!\left[
    \hat{p}_{y_n}^{(n)}\right]\right)^{\gamma(t)},\\
    \gamma(t)
    &= \begin{cases}
    0, & t < T_w, \\
    \frac{\gamma_{\max}}{2}
    \left(1+\cos\left(\pi
    \frac{t - T_w}{T - T_w}\right)\right), & t \ge T_w,
    \end{cases}\\
    \gamma_{\max} &= 5.0 .
\end{aligned}
\end{equation}
This is the effective EFL exponent. The implementation also logs a separate
topology smoothing parameter for the sparse structural controller. That logged
topology value remains at its base value during dense warmup and should not be
read as the effective focal exponent in Eq.~\ref{eq:focal_schedule}.
For KL regularization, we remove the true class evidence from the penalty by defining
$\tilde{\alpha}_c^{(n)}=y_c^{(n)}+(1-y_c^{(n)})\alpha_c^{(n)}$. The adjusted KL term is:
\begin{equation}\label{eq:kl_loss_main}
    L_{\text{KL}}^{(n)}
    =
    \KL\!\left(\Dir(\tilde{\bm{\alpha}}^{(n)})\parallel\Dir(\mathbf{1})\right).
\end{equation}
The annealing coefficient is $\mu(t)=\min(1.0,t/t_{\text{anneal}})$ with $t_{\text{anneal}}=10$. In the reported symmetric setting, the KL regularizer treats all classes uniformly, equivalent to multiplying $L_{\text{KL}}^{(n)}$ by $1.0$. The asymmetric KL diagnostic variant instead uses a capped class dependent multiplier:
\begin{equation}\label{eq:asym_multiplier}
    \Lambda_{\text{asym}}^{(n)}
    =
    \min(\bar{\omega}_{y_n},\Lambda_{\max}),
    \qquad \Lambda_{\max}=10.0 .
\end{equation}
The reported model applies $\bar{\omega}_{y_n}$ to the supervised CE term. The symmetric KL penalty has no additional class dependent multiplier. The asymmetric multiplier is a bounded diagnostic intervention. It is not part of the reported DST-EDL configuration.
The resulting training objective is
\begin{equation}\label{eq:total_loss}
    L_{\text{total}}=L_{\text{EFL}}(\bm{e},\bm{y},t).
\end{equation}

\paragraph{Controlled imbalance modulation.}
For each sample $n$, the stop gradient focal factor $\mathcal{F}_n(t)$ reweights the Dirichlet CE/KL gradients with respect to the evidence parameters. It does not introduce an additional derivative path through $\hat{p}_{y_n}^{(n)}$. The supervised class weight is bounded by $\bar{\omega}_{y_n}\le\omega_{\max}$. In the reported symmetric KL setting, the KL term has no additional class dependent multiplier. In the diagnostic asymmetric variant, the optional multiplier is bounded by $\Lambda_{\max}$. Focal modulation changes sample weights. It does not introduce an additional confidence gradient feedback loop through $\hat{p}_{y_n}$.

\paragraph{Training Warmup.}
We use a short learning rate and loss scale warmup for numerical stability. Its schedule is an implementation detail and is given with the full training algorithm in Appendix~\ref{app:algorithm}.


\subsection{Post hoc Calibration and Decision Thresholding}
Post hoc calibration is applied to logits before the Softplus evidence map. This changes predictive probabilities and Dirichlet strength. We treat it as evidential post hoc calibration. It is not ordinary softmax temperature scaling. Calibration is applied after training. It does not alter the learned sparse topology. The reported runs use bias temperature calibration with prior shift correction fitted on the calibration partition. Appendix~\ref{app:algorithm} gives the formula.

Operating thresholds are selected by sweeping validation probabilities $\hat p_c=\alpha_c/S$. Threshold selection changes the decision rule for Sensitivity and Specificity. It does not recalibrate the evidence values. The completed repeated seed ablations show different responses for sensitivity, specificity, ranking, and calibration. We report calibration and ranking metrics separately. We do not treat any single post hoc operating point as uniformly beneficial under extreme prevalence.

\paragraph{Scope of Claims.}
The validity checks above establish a well defined structural update. They establish mask feasibility and local first order saliency under a proper Dirichlet evidence state. They do not establish convergence, optimal topology discovery, Bayesian epistemic correctness, performance superiority over other sparse training rules, or realized sparse kernel acceleration. Empirical claims are limited to the reported ISIC 2024 case study. Broader validation protocols are treated as supporting evidence only when explicitly reported.

\section{Experimental Setup}\label{sec:experiment}

\subsection{Evaluation Benchmark}
The main empirical claims are centered on ISIC 2024. Appendix~\ref{app:external} reports a completed CIFAR-100-LT 1:100 diagnostic. Additional imbalance profiles, backbones, and realized sparse kernel profiling remain future validation targets unless explicitly reported. The current evidence supports feasibility and stability under the reported case study. It does not support a domain general SOTA claim.

\paragraph{ISIC 2024 Case Study.} We instantiate DST-EDL on the ISIC 2024 Challenge dataset~\cite{isic2024} (3D-TBP images with approximately 0.15\% minority prevalence). To prevent data leakage, we use a stratified 70/5/5/20 train/validation/calibration/test split grouped by patient\_id with split seed 42. The resulting partitions contain 5,775 training images, 21,930 validation images, 24,235 calibration images, and 65,931 held out test images. The test partition contains 77 malignant and 65,854 benign samples. The validation partition is used for decision threshold sweeps, the calibration partition for temperature and bias calibration, and the held out test partition for final reporting. Images are resized to $224\times224$ and normalized using ImageNet-1K statistics.

\paragraph{Metrics.} Default thresholds are unreliable under extreme prevalence skew. We report validation selected balanced operating points together with threshold independent ranking and calibration metrics. The main comparison reports Balanced Accuracy, pAUC$_{0.80}$, Macro AUROC, PR AUC, adaptive ECE, and valid block 2:4 feasibility. AURC is reported in the seed level and component diagnostic tables. Sensitivity and Specificity are reported when available for operating point diagnostics. Lower AURC indicates better risk ordering. Lower ECE indicates better calibration. AURC is computed from the empirical risk coverage curve using the same unnormalized scale across methods.

\subsection{Baselines and Backbones}
\paragraph{Baselines.} The reported ISIC tables include dense evidential baselines with completed metrics. These are Dense EDL, Fisher EDL, Flexible EDL, and R-EDL. The tables also include static and dynamic sparse 2:4 baselines. These are Static 2:4 EDL and RigL style 2:4. One DST-EDL model is reported. The extended ISIC suite also includes three seed softmax long tailed baselines: standard CE, focal loss, class balanced CE, balanced softmax, LDAM-DRW, cRT, and MiSLAS. These softmax baselines are reported as an appendix diagnostic. They are not used as the main comparison. They do not share the evidential output space or the exact 2:4 sparse topology mechanism. The completed CIFAR-100-LT 1:100 runs are a separate diagnostic. They compare Standard CE, Dense EDL, and the same reported DST-EDL model. They are not used as evidence for domain general superiority.

\paragraph{Backbone Instantiation and Wrapping.}
The reported case study uses a ResNet-18 instantiation. Inside each residual block, the $3 \times 3$ convolutional layers are replaced with sparse convolutional wrappers. These wrappers apply masks generated by latent scores. The first input convolution is kept dense. The structural proxy pass computes the risk ratio gradient for the pruner. It also computes KL to uniform gradient saliency for the regrower. Appendix~\ref{app:seed_level} documents additional sparse diagnostics. These diagnostics are not separate claimed contributions.

\paragraph{Implementation Details.} Models are trained using 16 bit mixed precision (AMP)~\cite{micikevicius2018mixed}. Evidential losses are evaluated in FP32. The reported ISIC models use AdamW with cosine annealing over 40 epochs. They use a one epoch learning rate warmup and the same validation/calibration/test protocol. For ISIC, the training partition subsamples the majority class to approximately 1:20. The evidential loss uses square root class frequency weights. Operating points are selected by validation threshold sweeps. The reported ECE column is adaptive ECE with 15 equal frequency bins. Equal width ECE is computed only as an auxiliary diagnostic.

\section{Results on the ISIC 2024 Case Study}\label{sec:results}

\subsection{Repeated Seed Evidential Comparison}
Table~\ref{tab:main_evidential_comparison} compares DST-EDL with dense evidential baselines. All rows use repeated seeds and the same split and threshold selection protocol. DST-EDL obtains the highest mean pAUC$_{0.80}$ and Macro AUROC in this ISIC 2024 comparison. It also preserves exact valid block 2:4 sparsity.

DST-EDL is not uniformly better than the dense evidential baselines. Its main advantage is high sensitivity ranking under exact structured sparsity. Mean pAUC$_{0.80}$ improves over Dense EDL ($0.1168$ vs. $0.0885$). Mean Macro AUROC improves from $0.8663$ to $0.8891$. Dense EDL remains slightly stronger in balanced accuracy and calibration. Several evidential baselines obtain higher PR AUC. The result indicates a ranking, specificity, and calibration tradeoff.

\begin{table}[t]
\centering
\caption{Compact ISIC 2024 repeated seed comparison under validation selected balanced operating points. Values are mean $\pm$ standard deviation over seeds 42, 123, and 456. Lower ECE is better.}
\label{tab:main_evidential_comparison}
\scriptsize
\renewcommand{\arraystretch}{1.03}
\setlength{\tabcolsep}{1.4pt}
\begin{adjustbox}{max width=\columnwidth}
\begin{tabular}{@{}l C C C C C@{}}
\toprule
\textbf{Model} & \textbf{Bal.} $\uparrow$ & \textbf{pAUC$_{0.80}$} $\uparrow$ & \textbf{Macro AUROC} $\uparrow$ & \textbf{PR AUC} $\uparrow$ & \textbf{ECE} $\downarrow$ \\
\midrule
Dense EDL & $\mathbf{0.8221\pm0.0195}$ & $0.0885\pm0.0124$ & $0.8663\pm0.0102$ & $0.0413\pm0.0085$ & $0.0524\pm0.0013$ \\
Fisher EDL & $0.8080\pm0.0152$ & $0.0842\pm0.0111$ & $0.8604\pm0.0163$ & $0.0414\pm0.0104$ & $\mathbf{0.0519\pm0.0009}$ \\
Flexible EDL & $0.8059\pm0.0116$ & $0.0847\pm0.0022$ & $0.8595\pm0.0078$ & $0.0373\pm0.0064$ & $0.0525\pm0.0013$ \\
R-EDL & $0.8034\pm0.0357$ & $0.0743\pm0.0175$ & $0.8479\pm0.0241$ & $\mathbf{0.0422\pm0.0051}$ & $0.0523\pm0.0012$ \\
\midrule
\textbf{DST-EDL} & $0.8110\pm0.0050$ & $\mathbf{0.1168\pm0.0070}$ & $\mathbf{0.8891\pm0.0080}$ & $0.0355\pm0.0063$ & $0.0564\pm0.0009$ \\
\bottomrule
\end{tabular}
\end{adjustbox}
\end{table}

\subsection{Sparse 2:4 Baseline Comparison}
Table~\ref{tab:sparse_comparison} compares DST-EDL against sparse 2:4 baselines with completed three seed metrics.

\begin{table}[t]
\centering
\caption{Sparse 2:4 comparison over seeds 42, 123, and 456. Values are mean $\pm$ standard deviation. All sparse models maintain 50\% active density in wrapped layers. Lower ECE is better.}
\label{tab:sparse_comparison}
\scriptsize
\renewcommand{\arraystretch}{1.03}
\setlength{\tabcolsep}{1.4pt}
\begin{adjustbox}{max width=\columnwidth}
\begin{tabular}{@{}l C C C C C@{}}
\toprule
\textbf{Model} & \textbf{Bal.} $\uparrow$ & \textbf{pAUC$_{0.80}$} $\uparrow$ & \textbf{Macro AUROC} $\uparrow$ & \textbf{PR AUC} $\uparrow$ & \textbf{ECE} $\downarrow$ \\
\midrule
Static 2:4 EDL & $0.7956\pm0.0121$ & $0.0965\pm0.0109$ & $0.8663\pm0.0149$ & $\mathbf{0.0444\pm0.0099}$ & $\mathbf{0.0529\pm0.0006}$ \\
RigL style 2:4 & $0.7696\pm0.0608$ & $0.1096\pm0.0100$ & $0.8829\pm0.0170$ & $0.0373\pm0.0034$ & $0.0541\pm0.0004$ \\
\textbf{DST-EDL} & $\mathbf{0.8110\pm0.0050}$ & $\mathbf{0.1168\pm0.0070}$ & $\mathbf{0.8891\pm0.0080}$ & $0.0355\pm0.0063$ & $0.0564\pm0.0009$ \\
\bottomrule
\end{tabular}
\end{adjustbox}
\end{table}

DST-EDL improves over Static 2:4 EDL in balanced accuracy ($0.8110$ vs. $0.7956$), pAUC$_{0.80}$ ($0.1168$ vs. $0.0965$), and Macro AUROC ($0.8891$ vs. $0.8663$). The result is consistent with a high sensitivity ranking benefit from uncertainty guided topology adaptation over a fixed valid block sparse mask. DST-EDL also improves over RigL style 2:4 in pAUC$_{0.80}$ ($0.1168$ vs. $0.1096$), Macro AUROC ($0.8891$ vs. $0.8829$), and balanced accuracy ($0.8110$ vs. $0.7696$). Its operating point variance is lower. Static 2:4 remains stronger in PR AUC and ECE. The sparse comparison indicates a ranking and stability advantage. It does not support uniform superiority.

\paragraph{Sparse morphology diagnostics.}
All sparse models preserved 50.0\% active density and showed no detected structural collapse (see Appendix~\ref{app:seed_level}). For the current DST-EDL runs, all wrapped sparse layers satisfied exact valid block 2:4 feasibility, with 2,789,376/2,789,376 valid blocks and no detected representational collapse. This diagnostic rules out a trivial explanation in which performance differences arise from invalid masks or collapsed sparse layers.
Appendix~\ref{app:seed_level} also reports the structural cost interpretation:
the current claim is theoretical sparse compatibility, not realized hardware
acceleration.

\subsection{Seed Level Stability and Operating Point Tradeoffs}
Appendix~\ref{app:seed_level} reports seed level results. Across seeds 42, 123, and 456, DST-EDL obtains mean balanced accuracy $0.8110$, sensitivity $0.7835$, specificity $0.8386$, pAUC$_{0.80}=0.1168$, Macro AUROC $=0.8891$, PR AUC $=0.0355$, AURC $=0.00023$, and ECE $=0.0564$.
DST-EDL is more stable than the RigL style sparse baseline at the balanced operating point. Its balanced accuracy, sensitivity, and specificity standard deviations are $0.0050$, $0.0418$, and $0.0399$. The corresponding RigL style 2:4 values are $0.0608$, $0.2088$, and $0.1251$.

\subsection{Component Attribution Diagnostics}
\label{subsec:ablation}
Completed three seed component diagnostics suggest that the KL uniform regrower
is the most important structural component: removing it causes the largest drop
in pAUC$_{0.80}$, Macro AUROC, PR AUC, AURC, and balanced accuracy. Removing the
pruner produces a smaller consistent degradation. Asymmetric KL mostly
shifts the operating point toward higher specificity. Removing EFL improves ECE
and specificity. It sharply lowers sensitivity. Removing anti crystallization
slightly increases pAUC. It makes the operating point less stable. These
diagnostic variants are not treated as separate reported DST-EDL models. They
attribute the components that support the single reported model. Full tables are
provided in Appendix~\ref{app:ablation}.

\section{Limitations and Broader Impact}\label{sec:limitations}
DST-EDL should not be interpreted as a universally dominant sparse classifier. On ISIC 2024 it obtains the highest mean pAUC$_{0.80}$ and Macro AUROC among the compared methods. It does not dominate all metrics. Dense EDL remains stronger in balanced accuracy, specificity, and calibration. Static 2:4 EDL obtains higher PR AUC and lower ECE. DST-EDL trades some specificity for higher sensitivity and high TPR ranking. At the balanced operating point, DST-EDL reaches mean sensitivity $0.7835$ and specificity $0.8386$. The reported Dense EDL and Static 2:4 specificities are $0.9158$ and $0.9140$. The current evidence is centered on ISIC 2024. Appendix~\ref{app:external} reports the completed CIFAR-100-LT 1:100 diagnostic. That diagnostic preserves exact 2:4 structure and improves ECE relative to dense CIFAR baselines. It is unfavorable on classification and ranking metrics relative to Standard CE and Dense EDL. It is reported as a limitation and stress test.

The ablations further narrow the claim. They are component diagnostics around the reported method. They are not competing full methods. The completed three seed variants show a central role for dormant link KL uniform regrowth in ranking and risk ordering. EFL trades calibration and specificity for rare event sensitivity. Anti crystallization acts mainly as a stability component. Calibration remains difficult under extreme prevalence. Exact valid block 2:4 feasibility does not imply realized acceleration without sparse Tensor Core compatible kernels and compatible layouts. In high stakes use, such models should support human decision making. They should not act as autonomous diagnostic systems. Dataset bias across skin types remains a concern.

\section{Conclusion}\label{sec:conclusion}
We introduced DST-EDL, a conservative uncertainty to topology sparse evidential framework for extremely imbalanced classification. The reported model combines topology caching, softened signed first order pruning, KL uniform regrowth, symmetric KL regularization, and evidential focal learning. It enforces exact valid block 2:4 sparsity. Its structural controller is supported by checks for Dirichlet signal validity, proxy batch risk ratio concentration, local first order saliency, gradient isolation, bounded latent scores, and exact valid block projection. On ISIC 2024, DST-EDL achieves the highest mean pAUC$_{0.80}=0.1168$ and Macro AUROC$=0.8891$ among the compared dense evidential and sparse 2:4 baselines. It preserves exact structured sparsity. It shows no detected representational collapse. The supported conclusion is narrow. DST-EDL achieves the highest mean high sensitivity ranking metrics in this comparison under exact 2:4 structured sparsity. It is not the best calibrated model. It is not the highest PR AUC model. Completed three seed component diagnostics identify KL uniform regrowth as the most critical structural component. They identify EFL and anti crystallization as tradeoff components. Broader validity requires additional datasets, backbones, ablations, and hardware aware sparse kernels.


\begin{thebibliography}{10}

\bibitem{guo2017calibration}
C.~Guo, G.~Pleiss, Y.~Sun, and K.~Q. Weinberger.
\newblock On calibration of modern neural networks.
\newblock In {\em International Conference on Machine Learning (ICML)}, 2017.

\bibitem{sensoy2018evidential}
M.~Sensoy, L.~Kaplan, and M.~Kandemir.
\newblock Evidential deep learning to quantify classification uncertainty.
\newblock In {\em Advances in Neural Information Processing Systems (NeurIPS)}, 2018.

\bibitem{malinin2018predictive}
A.~Malinin and M.~Gales.
\newblock Predictive uncertainty estimation via prior networks.
\newblock In {\em Advances in Neural Information Processing Systems (NeurIPS)}, 2018.

\bibitem{evci2020rigging}
U.~Evci, T.~Gale, J.~Menick, P.~S. Castro, and E.~Elsen.
\newblock Rigging the lottery: Making all tickets winners.
\newblock In {\em International Conference on Machine Learning (ICML)}, 2020.

\bibitem{cui2019class}
Y.~Cui, M.~Jia, T.-Y. Lin, Y.~Song, and S.~Belongie.
\newblock Class balanced loss based on effective number of samples.
\newblock In {\em IEEE/CVF Conference on Computer Vision and Pattern Recognition (CVPR)}, 2019.

\bibitem{cao2019learning}
K.~Cao, C.~Wei, A.~Gaidon, N.~Arechiga, and T.~Ma.
\newblock Learning imbalanced datasets with label distribution aware margin loss.
\newblock In {\em Advances in Neural Information Processing Systems (NeurIPS)}, 2019.

\bibitem{menon2021long}
A.~K. Menon, S.~Jayasumana, A.~S. Rawat, H.~Jain, A.~Veit, and S.~Kumar.
\newblock Long tail learning via logit adjustment.
\newblock In {\em International Conference on Learning Representations (ICLR)}, 2021.

\bibitem{ren2020balanced}
J.~Ren, C.~Yu, S.~Sheng, X.~Ma, H.~Zhao, S.~Yi, and H.~Li.
\newblock Balanced meta softmax for long tailed visual recognition.
\newblock In {\em Advances in Neural Information Processing Systems (NeurIPS)}, 2020.

\bibitem{kang2019decoupling}
B.~Kang, S.~Xie, M.~Rohrbach, Z.~Yan, A.~Gordo, J.~Feng, and Y.~Kalantidis.
\newblock Decoupling representation and classifier for long tailed recognition.
\newblock In {\em International Conference on Learning Representations (ICLR)}, 2020.

\bibitem{zhong2021improving}
Z.~Zhong, J.~Cui, S.~Liu, and J.~Jia.
\newblock Improving calibration for long tailed recognition.
\newblock In {\em IEEE/CVF Conference on Computer Vision and Pattern Recognition (CVPR)}, 2021.

\bibitem{zhou2020bbn}
B.~Zhou, Q.~Cui, X.-S. Wei, and Z.-M. Chen.
\newblock BBN: Bilateral branch network with cumulative learning for long tailed visual recognition.
\newblock In {\em IEEE/CVF Conference on Computer Vision and Pattern Recognition (CVPR)}, 2020.

\bibitem{cui2021parametric}
J.~Cui, Z.~Zhong, S.~Liu, B.~Yu, and J.~Jia.
\newblock Parametric contrastive learning.
\newblock In {\em IEEE/CVF International Conference on Computer Vision (ICCV)}, 2021.

\bibitem{deng2023fisher}
D.~Deng, G.~Chen, Y.~Yu, F.~Liu, and P.-A. Heng.
\newblock Uncertainty estimation by Fisher information based evidential deep learning.
\newblock In {\em International Conference on Machine Learning (ICML)}, PMLR, 2023.

\bibitem{chen2023redl}
M.~Chen, J.~Gao, and C.~Xu.
\newblock R-EDL: Relaxing nonessential settings of evidential deep learning.
\newblock In {\em International Conference on Learning Representations (ICLR)}, 2024.

\bibitem{chen2024reedl}
M.~Chen, J.~Gao, and C.~Xu.
\newblock Revisiting essential and nonessential settings of evidential deep learning.
\newblock {\em arXiv preprint arXiv:2410.00393}, 2024.

\bibitem{yoon2025flexible}
T.~Yoon and H.~Kim.
\newblock Uncertainty estimation by flexible evidential deep learning.
\newblock {\em arXiv preprint arXiv:2510.18322}, 2025.

\bibitem{yu2024anedl}
Y.~Yu, D.~Deng, F.~Liu, Y.~Jin, Q.~Dou, G.~Chen, and P.-A. Heng.
\newblock ANEDL: Adaptive negative evidential deep learning for open set semi supervised learning.
\newblock In {\em AAAI Conference on Artificial Intelligence}, 2024.

\bibitem{wu2024evidence}
Y.~Wu, B.~Shi, B.~Dong, Q.~Zheng, and H.~Wei.
\newblock The evidence contraction issue in deep evidential regression: Discussion and solution.
\newblock In {\em AAAI Conference on Artificial Intelligence}, 2024.

\bibitem{shen2024mirage}
M.~Shen, J.~J. Ryu, S.~Ghosh, Y.~Bu, P.~Sattigeri, S.~Das, and G.~W. Wornell.
\newblock Are uncertainty quantification capabilities of evidential deep learning a mirage?
\newblock In {\em Advances in Neural Information Processing Systems (NeurIPS)}, 2024.

\bibitem{geifman2017selective}
Y.~Geifman and R.~El-Yaniv.
\newblock Selective classification for deep neural networks.
\newblock In {\em Advances in Neural Information Processing Systems (NeurIPS)}, 2017.

\bibitem{geifman2019selectivenet}
Y.~Geifman and R.~El-Yaniv.
\newblock SelectiveNet: A deep neural network with an integrated reject option.
\newblock In {\em International Conference on Machine Learning (ICML)}, 2019.

\bibitem{ding2020revisiting}
Y.~Ding, J.~Liu, J.~Xiong, and Y.~Shi.
\newblock Revisiting the evaluation of uncertainty estimation and its application to explore model complexity uncertainty tradeoff.
\newblock In {\em IEEE/CVF Conference on Computer Vision and Pattern Recognition Workshops (CVPRW)}, 2020.

\bibitem{traub2024overcoming}
J.~Traub, T.~J. Bungert, C.~T. L{\"u}th, M.~Baumgartner, K.~H. Maier-Hein, L.~Maier-Hein, and P.~F. Jaeger.
\newblock Overcoming common flaws in the evaluation of selective classification systems.
\newblock In {\em Advances in Neural Information Processing Systems (NeurIPS)}, 2024.

\bibitem{dettmers2019sparse}
T.~Dettmers and L.~Zettlemoyer.
\newblock Sparse networks from scratch: Faster training without losing performance.
\newblock In {\em ICLR Workshop}, 2019.

\bibitem{lee2019snip}
N.~Lee, T.~Ajanthan, and P.~H.~S. Torr.
\newblock SNIP: Single shot network pruning based on connection sensitivity.
\newblock In {\em International Conference on Learning Representations (ICLR)}, 2019.

\bibitem{li2025pffdst}
L.~Li, H.~Yang, J.~Guo, H.~Yu, M.~Qin, and T.~Zhang.
\newblock An efficient and accurate dynamic sparse training framework based on parameter freezing.
\newblock In {\em AAAI Conference on Artificial Intelligence}, 2025.

\bibitem{lasby2023srigl}
M.~Lasby, A.~Golubeva, U.~Evci, M.~Nica, and Y.~Ioannou.
\newblock Dynamic sparse training with structured sparsity.
\newblock In {\em International Conference on Learning Representations (ICLR)}, 2024.

\bibitem{mishra2021accelerating}
Asit Mishra, Jorge Albericio Latorre, Jeff Pool, Darko Stosic, Dusan Stosic, Ganesh Venkatesh, Chong Yu, and Paulius Micikevicius.
\newblock Accelerating sparse deep neural networks.
\newblock {\em arXiv preprint arXiv:2104.08378}, 2021.

\bibitem{hu2024accelerating}
Y.~Hu, K.~Zhao, W.~Huang, J.~Chen, and J.~Zhu.
\newblock Accelerating transformer pre-training with 2:4 sparsity.
\newblock {\em arXiv preprint arXiv:2404.01847}, 2024.

\bibitem{nvidiacusparselt}
NVIDIA.
\newblock cuSPARSELt: A high-performance CUDA library for sparse matrix-matrix multiplication.
\newblock \url{https://docs.nvidia.com/cuda/cusparselt/}, accessed 2026.

\bibitem{pytorchsemistructured}
PyTorch.
\newblock Semi-structured sparsity and \texttt{to\_sparse\_semi\_structured}.
\newblock \url{https://docs.pytorch.org/docs/2.12/sparse.html}, accessed 2026.

\bibitem{mocanu2018scalable}
D.~C. Mocanu, E.~Mocanu, P.~Stone, P.~H. Nguyen, M.~Gibescu, and A.~Liotta.
\newblock Scalable training of artificial neural networks with adaptive sparse connectivity inspired by network science.
\newblock {\em Nature Communications}, 9(1):2383, 2018.

\bibitem{han2016deep}
S.~Han, H.~Mao, and W.~J. Dally.
\newblock Deep compression: Compressing deep neural networks with pruning, trained quantization and Huffman coding.
\newblock In {\em International Conference on Learning Representations (ICLR)}, 2016.

\bibitem{frankle2019lottery}
J.~Frankle and M.~Carbin.
\newblock The lottery ticket hypothesis: Finding sparse, trainable neural networks.
\newblock In {\em International Conference on Learning Representations (ICLR)}, 2019.

\bibitem{liu2019large}
Z.~Liu, Z.~Miao, X.~Zhan, J.~Wang, B.~Gong, and S.~X. Yu.
\newblock Large scale long tailed recognition in an open world.
\newblock In {\em IEEE/CVF Conference on Computer Vision and Pattern Recognition (CVPR)}, 2019.


\bibitem{isic2024}
International Skin Imaging Collaboration (ISIC).
\newblock ISIC 2024 -- Skin Cancer Detection with 3D-TBP.
\newblock {\em Kaggle Competition}, 2024.

\bibitem{jaeger2023call}
P.~F. Jaeger, C.~T. L{\"u}th, L.~Klein, and T.~J. Bungert.
\newblock A call to reflect on evaluation practices for failure detection in image classification.
\newblock In {\em International Conference on Learning Representations (ICLR)}, 2023.

\bibitem{zhou2023domain}
M.~Zhou, A.~Jamzad, J.~Izard, A.~Menard, R.~Siemens, and P.~Mousavi.
\newblock Domain transfer through image to image translation for uncertainty aware prostate cancer classification.
\newblock {\em arXiv preprint arXiv:2307.00479}, 2023.

\bibitem{micikevicius2018mixed}
P.~Micikevicius, S.~Narang, J.~Alben, G.~Diamos, E.~Elsen, D.~Garcia, B.~Ginsburg, M.~Houston, O.~Kuchaiev, G.~Venkatesh, and H.~Wu.
\newblock Mixed precision training.
\newblock In {\em International Conference on Learning Representations (ICLR)}, 2018.

\bibitem{loshchilov2019decoupled}
I.~Loshchilov and F.~Hutter.
\newblock Decoupled weight decay regularization.
\newblock In {\em International Conference on Learning Representations (ICLR)}, 2019.

\bibitem{kingma2015adam}
D.~P. Kingma and J.~Ba.
\newblock Adam: A method for stochastic optimization.
\newblock In {\em International Conference on Learning Representations (ICLR)}, 2015.

\bibitem{charpentier2020posterior}
B.~Charpentier, D.~Z{\"u}gner, and S.~G{\"u}nnemann.
\newblock Posterior Network: Uncertainty estimation without OOD samples via density based pseudo counts.
\newblock In {\em Advances in Neural Information Processing Systems (NeurIPS)}, 2020.

\end{thebibliography}

\appendix

\section*{Appendix Roadmap}
The appendix follows the order most useful for reproducibility and review.
Appendix~A gives the experimental setting, split rules, hyperparameters, and
release plan. Appendix~B provides the training loop and implementation
configuration. Appendix~C reports seed level ISIC results and sparse morphology
checks. Appendix~D contains component attribution diagnostics. Appendix~E gives
well posedness and algorithmic validity checks. Appendix~F clarifies the novelty
boundary against adjacent methods. Appendix~G reports additional softmax and
CIFAR diagnostics together with future validation protocols.

\section{Reproducibility Details}
\label{app:reproducibility}

\subsection*{1. Experimental Setting and Hyperparameters}
\begin{itemize}
    \item \textbf{Random Seeds:} Main ISIC comparisons use model seeds 42, 123, and 456 with split seed 42. Paper level component ablations include only variants that completed all three model seeds. Partial single seed or two seed ablations are excluded from the reported ablation table.
    \item \textbf{Preprocessing \& Dimensions:} Images are resized to $224 \times 224$ and normalized using ImageNet channel wise statistics: $\mu = [0.485, 0.456, 0.406]$ and $\sigma = [0.229, 0.224, 0.225]$. Data augmentation includes horizontal and vertical flips.
    \item \textbf{Optimizer and Learning Rate Schedule:} We employ the AdamW optimizer ($\beta_1=0.9$, $\beta_2=0.999$, weight decay $10^{-4}$ on active weights). The learning rate is initialized at $1\times10^{-6}$ and linearly warmed up to $4\times10^{-5}$ over the first epoch, followed by a cosine annealing schedule over 40 epochs.
\end{itemize}

\subsection*{2. Evaluation Protocols}
\begin{itemize}
    \item \textbf{Validation/Test Split Rules:} The ISIC 2024 dataset is partitioned into 70\% training, 5\% validation, 5\% calibration, and 20\% test sets using split seed 42. The resulting split contains 5,775 training images, 21,930 validation images, 24,235 calibration images, and 65,931 test images, with 77 malignant and 65,854 benign test cases. To guarantee zero patient level data leakage, splits are strictly grouped by patient ID using a stratified grouping split protocol. The test set is held out from training, calibration, model selection, and threshold selection, and is used only for final reporting after each model configuration is fixed.
    \item \textbf{Threshold Selection Protocol:} For ISIC, operating points are selected by performing threshold sweeps on the isolated validation partition to hit target clinical sensitivities ($\ge 80\%$). Temperature and bias parameters are fitted on the separate calibration partition. Validation derived thresholds are frozen before final test evaluation.
    \item \textbf{ECE Binning and pAUC:} The reported ECE is adaptive ECE with 15 equal frequency bins. Equal width ECE is also computed as an auxiliary diagnostic. Partial AUROC is computed strictly within the True Positive Rate interval $[0.80, 1.0]$.
\end{itemize}

\subsection*{3. Source Code and Data Releases}
\begin{itemize}
    \item \textbf{Code Availability:} Full PyTorch code containing wrapper layers, evidential loss definitions, and evaluation scripts is packaged in a supplementary repository and will be fully open sourced on GitHub upon publication.
\end{itemize}

\section{DST-EDL Algorithm and Implementation Details}
\label{app:algorithm}
This appendix gives the procedural details that are referenced from the
optimization section. It contains the warmup rule, the post hoc evidential
calibration formula, the full training loop, and the implementation
configuration used for the reported ISIC 2024 case study.

\paragraph{Learning rate and loss scale warmup.}
During the first epoch, the learning rate is linearly warmed from $\eta_{\text{init}}=10^{-6}$ to $\eta_{\text{base}}=4\times10^{-5}$. An initial loss multiplier decays from $S_{\max}=4.0$ to $1.0$:
\begin{equation}\label{eq:warmup_main}
\begin{aligned}
    \eta_s
    &=
    \eta_{\text{init}}
    +(\eta_{\text{base}}-\eta_{\text{init}})
    \frac{s}{S_{\text{warmup}}},\\
    S_{\text{loss}}(s)
    &=
    S_{\max}-(S_{\max}-1)
    \frac{s}{S_{\text{warmup}}}.
\end{aligned}
\end{equation}
The scaled objective $L_{\text{scaled}}=S_{\text{loss}}(s)L_{\text{total}}$ is used during warmup and $L_{\text{total}}$ afterward.

\paragraph{Bias corrected evidential calibration.}
After training, logits are corrected before the Softplus evidence map by
\begin{equation}
\begin{aligned}
    z^{\mathrm{cal}}_c
    &=
    \frac{z^{\mathrm{base}}_c+\Delta^\pi_c}{T_{\mathrm{cal}}}
    + b_{\mathrm{res},c},\\
    \Delta^\pi_c
    &=
    \log(\pi_{\mathrm{true},c}+\epsilon)
    -\log(\pi_{\mathrm{train},c}+\epsilon).
\end{aligned}
\end{equation}
Temperature only calibration sets $\bm{b}_{\mathrm{res}}=\bm{0}$. The reported bias temperature mode learns both $T_{\mathrm{cal}}$ and $\bm{b}_{\mathrm{res}}$ on the calibration partition.

\begin{table}[!t]
\centering
\caption{DST-EDL training and optimization loop.}
\label{alg:dst_edl}
\scriptsize
\renewcommand{\arraystretch}{1.04}
\begin{adjustbox}{max width=\columnwidth}
\begin{tabular}{@{}r p{0.82\columnwidth}@{}}
\toprule
1 & Initialize data $\mathcal{D}$, weights $\bm{W}^{(0)}$, latent scores $\bm{V}^{(0)}$, epochs $T$, and warmup $T_w$. \\
2 & For each epoch $t=1,\ldots,T$, freeze channel maps when $t=T_w$. \\
3 & If $t\ge T_w$, sample a proxy batch $(\bm{X},\bm{y})\sim\mathcal{D}$. \\
4 & Compute $R_{\mathcal B}$ and $L_{\text{grow}}$, then store pruner and regrower gradients. \\
5 & Compute $C_{ij}$ by Eq.~\ref{eq:C_ij} and $G_{ij}$ by Eq.~\ref{eq:G_ij}. \\
6 & Add small exploration noise when $G_{ij}$ is nearly flat. \\
7 & Update latent scores $\bm{V}$ by Eq.~\ref{eq:score_update}. \\
8 & Cache 2:4 masks $\bm{M}^{(l)}$ using Eq.~\ref{eq:mask}. During warmup set masks to dense. \\
9 & For each mini batch, compute $\eta_s$ and $S_{\text{loss}}(s)$ by Eq.~\ref{eq:warmup_main}. \\
10 & Forward with $\bm{W}_{\text{eff}}=\bm{W}\odot\bm{M}$, set $\alpha_c=\ln(1+e^{z_c})+1$, and compute uncertainties. \\
11 & Compute $L_{\text{scaled}}=S_{\text{loss}}(s)L_{\text{EFL}}$ using Eq.~\ref{eq:efl} and update active weights. \\
\bottomrule
\end{tabular}
\end{adjustbox}
\end{table}

\begin{table}[!t]
\centering
\caption{Training and calibration configuration for the ISIC 2024 case study.}
\label{tab:config}
\scriptsize
\renewcommand{\arraystretch}{0.96}
\begin{adjustbox}{max width=\columnwidth}
\begin{tabular}{@{}p{0.30\columnwidth} p{0.24\columnwidth} p{0.36\columnwidth}@{}}
\toprule
\textbf{Component} & \textbf{Value} & \textbf{Role} \\
\midrule
Backbone & ResNet-18 & ImageNet-1K pretrained \\
Classes ($K$) & 2 & Benign / malignant \\
Dense warmup mask & $\bm{1}$ & No pruning before $T_w$ \\
Warmup EFL exponent & $0.0$ & Effective focal exponent during dense warmup \\
2:4 mask & top two of four & Project latent scores $\bm{V}$ \\
Topology refresh & Cached once/epoch & Up to 4 train mini batches \\
Post hoc calibration & Bias temperature & Fit $T_{\text{cal}}$ and residual bias \\
Optimizer & AdamW & $\beta_1=0.9$, $\beta_2=0.999$ \\
Learning rate & $4{\times}10^{-5}$ & Linear warmup from $10^{-6}$ \\
Weight decay & $10^{-4}$ & Active weights \\
Gradient clipping & $1.0$ & Global $L_2$ norm \\
Batch size / epochs & 32 / 40 & ISIC repeated seed runs \\
Dense topology warmup & 12 epochs & 30\% of training budget \\
EFL $\gamma_{\text{base}}$ & $5.0$ & Cosine decay after warmup \\
$\lambda_{\text{KL}}$, $t_{\text{anneal}}$ & $0.1$, 10 & KL strength and annealing \\
$\omega_{\max}$ & $10.0$ & Cap class weights \\
\bottomrule
\end{tabular}
\end{adjustbox}
\end{table}

\begin{table}[!t]
\centering
\caption{Structural controller hyperparameters for reported DST-EDL runs.}
\label{tab:struct_config}
\scriptsize
\renewcommand{\arraystretch}{0.98}
\begin{adjustbox}{max width=\columnwidth}
\begin{tabular}{@{}p{0.34\columnwidth} p{0.23\columnwidth} p{0.35\columnwidth}@{}}
\toprule
\textbf{Parameter} & \textbf{Value} & \textbf{Role} \\
\midrule
$\beta_m$ & $0.95$ & Structural momentum \\
$\eta_{\text{struct}}$ & $0.02$ & Latent score step size \\
$\rho$ & $0.25$ & Dormant link attenuation \\
$V_{\max}$ & $5.0$ & Latent score clamp \\
Pruner & Soft signed first order & Eq.~\ref{eq:C_ij}, strength $0.5$ \\
Regrower & KL uniform & Eq.~\ref{eq:grow_objective} \\
Exploration scale & $0.0316$ & Anti crystallization near flat scores \\
\bottomrule
\end{tabular}
\end{adjustbox}
\end{table}

\FloatBarrier

\section{Seed Level Results and Sparse Morphology Diagnostics}
\label{app:seed_level}
This appendix provides the per seed evidence behind the aggregated ISIC 2024
numbers in the main paper, followed by the sparse morphology checks used to
verify exact valid block 2:4 feasibility. The paragraph is intentionally placed
before the seed level tables. This appendix has visible body text when
LaTeX moves one of the floats to the next column or page.

\begin{table}[!t]
\centering
\caption{Seed level DST-EDL operating points on ISIC 2024 for split seed 42. Lower ECE is better.}
\label{tab:seed_level}
\scriptsize
\renewcommand{\arraystretch}{1.05}
\setlength{\tabcolsep}{2.0pt}
\begin{adjustbox}{max width=\columnwidth}
\begin{tabular}{@{}l C C C C@{}}
\toprule
\textbf{Seed} & \textbf{Bal.} $\uparrow$ & \textbf{Sens.} $\uparrow$ & \textbf{Spec.} $\uparrow$ & \textbf{ECE} $\downarrow$ \\
\midrule
42 & $0.8155$ & $0.7662$ & $0.8649$ & $0.0555$ \\
123 & $0.8119$ & $0.8312$ & $0.7927$ & $0.0565$ \\
456 & $0.8057$ & $0.7532$ & $0.8581$ & $0.0573$ \\
\midrule
Mean $\pm$ Std & $0.8110\pm0.0050$ & $0.7835\pm0.0418$ & $0.8386\pm0.0399$ & $0.0564\pm0.0009$ \\
\bottomrule
\end{tabular}
\end{adjustbox}
\end{table}

\begin{table}[!t]
\centering
\caption{Seed level DST-EDL ranking and risk metrics on ISIC 2024. Lower AURC is better.}
\label{tab:seed_level_ranking}
\scriptsize
\renewcommand{\arraystretch}{1.05}
\setlength{\tabcolsep}{2.0pt}
\begin{adjustbox}{max width=\columnwidth}
\begin{tabular}{@{}l C C C C@{}}
\toprule
\textbf{Seed} & \textbf{pAUC$_{0.80}$} $\uparrow$ & \textbf{Macro AUROC} $\uparrow$ & \textbf{PR AUC} $\uparrow$ & \textbf{AURC} $\downarrow$ \\
\midrule
42 & $0.1244$ & $0.8982$ & $0.0421$ & $0.0002$ \\
123 & $0.1105$ & $0.8834$ & $0.0349$ & $0.0003$ \\
456 & $0.1155$ & $0.8857$ & $0.0296$ & $0.0002$ \\
\midrule
Mean $\pm$ Std & $0.1168\pm0.0070$ & $0.8891\pm0.0080$ & $0.0355\pm0.0063$ & $0.00023\pm0.00006$ \\
\bottomrule
\end{tabular}
\end{adjustbox}
\end{table}

Table~\ref{tab:structural_diagnostics} reports the corresponding structural
validity diagnostics for the sparse models.

\begin{table}[!t]
\centering
\caption{Structural diagnostics on ISIC 2024. Sparse models maintained 50.0\% active density with no detected structural collapse. The current DST-EDL runs satisfy all complete valid 2:4 blocks in the wrapped sparse layers.}
\label{tab:structural_diagnostics}
\footnotesize
\renewcommand{\arraystretch}{1.1}
\setlength{\tabcolsep}{2pt}
\begin{adjustbox}{max width=\columnwidth}
\begin{tabular}{@{}l c c c@{}}
\toprule
\textbf{Model} & \textbf{Density} & \textbf{Valid 2:4} & \textbf{Collapse} \\
\midrule
Static 2:4 EDL & 50.0\% & Exact 2:4 & No \\
RigL style 2:4 & 50.0\% & Exact 2:4 & No \\
\textbf{DST-EDL} & 50.0\% & Exact 2:4 & No \\
\bottomrule
\end{tabular}
\end{adjustbox}
\end{table}

\begin{table}[!t]
\centering
\caption{Structural cost interpretation. Theoretical sparse MAC reduction is reported only for wrapped sparse layers. Real speedup is not claimed. The current implementation uses masked PyTorch kernels and not sparse Tensor Core compatible kernels and layouts.}
\label{tab:structural_cost}
\scriptsize
\renewcommand{\arraystretch}{1.05}
\setlength{\tabcolsep}{1.6pt}
\begin{adjustbox}{max width=\columnwidth}
\begin{tabular}{@{}l c c c c@{}}
\toprule
\textbf{Model} & \textbf{Active} & \textbf{Valid 2:4} & \textbf{Sparse MAC} & \textbf{Speedup?} \\
\midrule
Dense EDL & 100\% & N/A & 0\% & No \\
Static 2:4 EDL & 50\% & Exact & 50\% wrapped & No \\
RigL style 2:4 & 50\% & Exact & 50\% wrapped & No \\
\textbf{DST-EDL} & 50\% & 2,789,376/2,789,376 & 50\% wrapped & No \\
\bottomrule
\end{tabular}
\end{adjustbox}
\end{table}

We report theoretical sparse compatibility. We do not report realized acceleration.
The current implementation uses masked PyTorch kernels. Real acceleration
requires sparse Tensor Core compatible kernels and layouts such as cuSPARSELt or
TensorRT sparse execution.

\FloatBarrier

\section{Component Attribution Diagnostics}
\label{app:ablation}
This appendix collects diagnostic component variants used for attribution.
They are not additional reported methods. Table~\ref{tab:ablation_components_full}
reports only variants that completed all three model seeds 42, 123, and 456
under split seed 42. Variants without three completed seeds are intentionally
omitted from the paper level ablation table.

\begin{table}[!t]
\centering
\caption{Three seed component operating point diagnostics on ISIC 2024. Only variants with all seeds 42, 123, and 456 are included. Lower ECE is better.}
\label{tab:ablation_components_full}
\scriptsize
\renewcommand{\arraystretch}{1.05}
\setlength{\tabcolsep}{1.8pt}
\begin{adjustbox}{max width=\columnwidth}
\begin{tabular}{@{}l C C C C@{}}
\toprule
\textbf{Variant} & \textbf{Bal.} $\uparrow$ & \textbf{Sens.} $\uparrow$ & \textbf{Spec.} $\uparrow$ & \textbf{ECE} $\downarrow$ \\
\midrule
\textbf{DST-EDL} & $\mathbf{0.8110\pm0.0050}$ & $\mathbf{0.7835\pm0.0418}$ & $0.8386\pm0.0399$ & $0.0564\pm0.0009$ \\
w/o pruner & $0.7946\pm0.0077$ & $0.7229\pm0.0492$ & $0.8662\pm0.0394$ & $0.0552\pm0.0007$ \\
w/o regrower & $0.7046\pm0.0527$ & $0.6753\pm0.0723$ & $0.7339\pm0.0520$ & $0.0556\pm0.0027$ \\
asymmetric KL & $0.8033\pm0.0116$ & $0.7229\pm0.0715$ & $0.8836\pm0.0485$ & $0.0556\pm0.0003$ \\
w/o EFL & $0.7500\pm0.0266$ & $0.5628\pm0.1050$ & $\mathbf{0.9372\pm0.0519}$ & $\mathbf{0.0544\pm0.0003}$ \\
w/o anti cryst. & $0.7837\pm0.0369$ & $0.6753\pm0.1498$ & $0.8920\pm0.0761$ & $0.0556\pm0.0006$ \\
\bottomrule
\end{tabular}
\end{adjustbox}
\end{table}

\begin{table}[!t]
\centering
\caption{Three seed component ranking and risk diagnostics on ISIC 2024. Lower AURC is better.}
\label{tab:ablation_ranking}
\scriptsize
\renewcommand{\arraystretch}{1.05}
\setlength{\tabcolsep}{1.8pt}
\begin{adjustbox}{max width=\columnwidth}
\begin{tabular}{@{}l C C C C@{}}
\toprule
\textbf{Variant} & \textbf{pAUC$_{0.80}$} $\uparrow$ & \textbf{Macro AUROC} $\uparrow$ & \textbf{PR AUC} $\uparrow$ & \textbf{AURC} $\downarrow$ \\
\midrule
\textbf{DST-EDL} & $0.1168\pm0.0070$ & $\mathbf{0.8891\pm0.0080}$ & $\mathbf{0.0355\pm0.0063}$ & $\mathbf{0.00023\pm0.00006}$ \\
w/o pruner & $0.1051\pm0.0057$ & $0.8720\pm0.0069$ & $0.0319\pm0.0029$ & $0.00027\pm0.00006$ \\
w/o regrower & $0.0640\pm0.0187$ & $0.7663\pm0.0629$ & $0.0077\pm0.0049$ & $0.00057\pm0.00021$ \\
asymmetric KL & $0.1127\pm0.0141$ & $0.8845\pm0.0159$ & $0.0348\pm0.0042$ & $0.00033\pm0.00006$ \\
w/o EFL & $0.0898\pm0.0010$ & $0.8546\pm0.0016$ & $0.0298\pm0.0038$ & $0.00030\pm0.00000$ \\
w/o anti cryst. & $\mathbf{0.1173\pm0.0098}$ & $0.8857\pm0.0109$ & $0.0352\pm0.0047$ & $0.00030\pm0.00010$ \\
\bottomrule
\end{tabular}
\end{adjustbox}
\end{table}

The three seed ablation gives a clearer attribution picture than the earlier
partial diagnostics. DST-EDL provides the best overall operating point
tradeoff. It is best on balanced accuracy, sensitivity, Macro AUROC, PR AUC,
and AURC. Its pAUC$_{0.80}$ is only slightly below the no anti crystallization
variant ($0.1168$ vs. $0.1173$). Removing the KL uniform regrower causes the
largest degradation. pAUC$_{0.80}$ decreases from $0.1168$ to $0.0640$.
Macro AUROC decreases from $0.8891$ to $0.7663$. PR AUC decreases from
$0.0355$ to $0.0077$. AURC worsens from $0.00023$ to $0.00057$. These results
identify dormant link regrowth as a core component. Removing the pruner causes a
smaller consistent reduction in balanced accuracy, sensitivity,
pAUC$_{0.80}$, and Macro AUROC. This indicates that pruning improves active
link selection. Regrowth remains the more critical component. Asymmetric KL
remains close to DST-EDL in ranking metrics. It shifts the operating point
toward higher specificity and lower sensitivity. This result favors the symmetric
KL choice for the reported high sensitivity setting. Removing EFL yields the
best ECE and highest specificity. Sensitivity drops sharply from $0.7835$ to
$0.5628$. This identifies EFL as a rare event sensitivity mechanism with a
calibration and specificity tradeoff. Removing anti crystallization slightly
improves pAUC$_{0.80}$. It produces much higher sensitivity variance. The
exploration term is interpreted as a stability oriented component.

\FloatBarrier

\section{Well Posedness and Algorithmic Validity Checks}
\label{app:proofs}

\begin{table}[!t]
\centering
\caption{Scope of formal claims made for DST-EDL.}
\label{tab:guarantees}
\small
\renewcommand{\arraystretch}{1.12}
\begin{tabular}{@{}p{0.47\columnwidth}p{0.43\columnwidth}@{}}
\toprule
\textbf{Established by design or proof} & \textbf{Not claimed} \\
\midrule
Positive Dirichlet parameters, bounded vacuity, and non negative ambiguity proxy &
Exact Bayesian epistemic correctness \\
Exact valid block 2:4 feasibility over eligible complete 4-blocks &
Global convergence or topology optimality \\
Bounded latent topology scores and deterministic top two projection &
Uniform superiority over dense or dynamic sparse baselines \\
Ordinary masked gradient isolation plus virtual probe structural saliency &
Validation improvement without empirical evidence \\
Local first order pruning/regrowth interpretations and proxy batch concentration &
Realized hardware speedup without compatible sparse kernels \\
\bottomrule
\end{tabular}
\end{table}

\subsection{Compact Dirichlet Identities}
For completeness, the Dirichlet density for $\bm{p}$ on the probability simplex is
\begin{equation}\label{eq:dirichlet_pdf_app}
\begin{aligned}
    \Dir(\bm{p}\mid\bm{\alpha})
    &=
    \frac{1}{B(\bm{\alpha})}
    \prod_{c=1}^{K}p_c^{\alpha_c - 1}\\
    &=
    \frac{\Gamma(S)}{\prod_{c=1}^{K}\Gamma(\alpha_c)}
    \prod_{c=1}^{K}p_c^{\alpha_c - 1}.
\end{aligned}
\end{equation}
The implementation uses the following identities from the main text:
\begin{equation}
\begin{aligned}
    \hat{p}_c&=\frac{\alpha_c}{S},
    &
    u_e&=\frac{K}{S},\\
    u_a&=\sum_{c=1}^{K}\frac{\alpha_c}{S}
    \left[\psi(S+1)-\psi(\alpha_c+1)\right].
\end{aligned}
\end{equation}
Under the Softplus evidence parameterization, $\alpha_c>1$ for finite logits. Therefore $S>K$. Hence $u_e\in(0,1)$. Its supremum $1$ is approached as all evidence tends to zero. Also, $S\ge \alpha_c$ and $\psi(\cdot)$ is monotone on positive inputs. Each entropy summand is non negative. Hence $u_a\ge0$.

\subsection{Structural Risk Ratio Chain Rule}
\label{app:r_chain}
The batch structural risk ratio $R_{\mathcal B}$ in Eq.~\ref{eq:batch_risk} is differentiable for $\epsilon>0$ under the Softplus evidence map when the current mask is held fixed during the proxy structural pass. Write $\partial_{ij}$ for the derivative with respect to $w_{\text{eff},ij}$. The weight derivative is the batch average of the per sample ratio derivatives:
\begin{equation}
\begin{aligned}
    \partial_{ij} R_{\mathcal B}
    &=
    \frac{1}{|\mathcal B|}\sum_{n\in\mathcal B}
    \left[
    \frac{1}{u_e+\epsilon}\partial_{ij}u_a
    -
    \frac{u_a}{(u_e+\epsilon)^2}
    \partial_{ij}u_e
    \right]^{(n)} .
\end{aligned}
\end{equation}
At the scalar level, $r(u_a,u_e)=u_a/(u_e+\epsilon)$ satisfies
\[
    \frac{\partial r}{\partial u_a}
    =
    \frac{1}{u_e+\epsilon}>0,
    \qquad
    \frac{\partial r}{\partial u_e}
    =
    -\frac{u_a}{(u_e+\epsilon)^2}\le 0,
\]
The proxy is monotone increasing in ambiguity and monotone decreasing in vacuity whenever $u_a\ge0$.
The vacuity derivative expands through the Dirichlet parameters as
\begin{equation}
    \frac{\partial u_e}{\partial w_{\text{eff},ij}}
    =
    \sum_{m=1}^{K}
    \frac{\partial u_e}{\partial \alpha_m}
    \frac{\partial \alpha_m}{\partial z_m}
    \frac{\partial z_m}{\partial w_{\text{eff},ij}},
    \qquad
    \frac{\partial u_e}{\partial \alpha_m}
    =
    -\frac{K}{S^2},
\end{equation}
where $\frac{\partial \alpha_m}{\partial z_m}=1/(1+\exp(-z_m))$ for $\alpha_m=\ln(1+\exp(z_m))+1$. For the entropy proxy, let $\psi_1(\cdot)$ denote the trigamma function and $\delta_{cm}$ the Kronecker delta. Then
\begin{equation}
\begin{aligned}
    \frac{\partial u_a}{\partial \alpha_m}
    ={}&
    \sum_{c=1}^{K}
    \frac{\delta_{cm}S-\alpha_c}{S^2}
    \left[\psi(S+1)-\psi(\alpha_c+1)\right]\\
    &+
    \sum_{c=1}^{K}
    \frac{\alpha_c}{S}
    \left[
    \psi_1(S+1)-\delta_{cm}\psi_1(\alpha_c+1)
    \right],
\end{aligned}
\end{equation}
This gives
\begin{equation}
    \frac{\partial u_a}{\partial w_{\text{eff},ij}}
    =
    \sum_{m=1}^{K}
    \frac{\partial u_a}{\partial \alpha_m}
    \frac{\partial \alpha_m}{\partial z_m}
    \frac{\partial z_m}{\partial w_{\text{eff},ij}} .
\end{equation}
This expansion provides the batch risk ratio derivative used by the signed pruning score in Eq.~\ref{eq:C_ij}. Throughout the structural proxy derivations, derivatives are taken with respect to the effective coordinate $w_{\text{eff},ij}$. When the mask is fixed and the coordinate is active, this corresponds to the ordinary weight derivative up to the fixed mask multiplier. The expression combines local vacuity and expected entropy derivatives through the ordinary network Jacobian $\partial z_m/\partial w_{\text{eff},ij}$. It does not rely on an unspecified pruning signal.

\subsection{Proxy Batch Concentration of the Risk Ratio Signal}
\label{app:proxy_concentration}
Fix the current weights and mask during a structural refresh, and let
$r(x)=u_a(x)/(u_e(x)+\epsilon)$ under the proxy sampling distribution.
For a $K$-class categorical variable, entropy is bounded above by
$\log K$, hence the Dirichlet expected categorical entropy satisfies
$0\le u_a(x)\le\log K$. Since $u_e(x)+\epsilon\ge\epsilon$, we have
\[
    0\le r(x)\le \frac{\log K}{\epsilon}.
\]
Let $\bar R=\mathbb{E}[r(x)]$ and
$R_{\mathcal B}=|\mathcal B|^{-1}\sum_{x\in\mathcal B}r(x)$ for an
i.i.d. proxy batch from this sampling distribution. Hoeffding's
inequality gives, for any $\delta\in(0,1)$, with probability at least
$1-\delta$,
\[
    |R_{\mathcal B}-\bar R|
    \le
    \frac{\log K}{\epsilon}
    \sqrt{\frac{\log(2/\delta)}{2|\mathcal B|}}.
\]
This result only controls the estimation noise of the scalar proxy under
the chosen proxy sampling distribution. It does not imply that the
resulting structural update improves validation risk, converges, or
dominates magnitude- or task gradient based sparse training.

\subsection{KL Regularization Note}
The reported DST-EDL configuration uses the symmetric KL penalty in Eq.~\ref{eq:kl_loss_main}: the KL multiplier is $1.0$ for every sample and class. The adjusted false evidence parameter $\tilde{\alpha}_c^{(n)}=y_c^{(n)}+(1-y_c^{(n)})\alpha_c^{(n)}$ defines the compact KL term. Its expanded form is the standard Dirichlet to Dirichlet KL with the true class concentration removed from the shrinkage penalty:
\begin{equation}\label{eq:kl_loss_expanded}
\begin{aligned}
    L_{\text{KL}}^{(n)}
    ={}& \ln \left(
    \frac{\Gamma\left(\sum_{c=1}^K\tilde{\alpha}_c^{(n)}\right)}
    {\Gamma(K)\prod_{c=1}^K\Gamma(\tilde{\alpha}_c^{(n)})}
    \right)\\
    &+\sum_{c=1}^K(\tilde{\alpha}_c^{(n)}-1)
    \left[
    \psi(\tilde{\alpha}_c^{(n)})
    -\psi\!\left(\sum_{c=1}^K\tilde{\alpha}_c^{(n)}\right)
    \right].
\end{aligned}
\end{equation}

\subsection{Validity of Evidential State}
Softplus is strictly positive for finite real valued logits. Hence $\alpha_c=e_c+1>1$ and $S=\sum_c\alpha_c>K$. The expected probabilities obey $\hat{p}_c=\alpha_c/S>0$ and $\sum_c\hat{p}_c=1$. Thus $\hat{\bm{p}}\in\Delta^K$. Since $S>K$, $u_e=K/S\in(0,1)$ for finite logits. Its limiting value $1$ is approached when all evidence values $e_c$ tend to zero. Finally, $S\ge\alpha_c$ for every class and $\psi(\cdot)$ is monotone on positive inputs. Therefore $\psi(S+1)-\psi(\alpha_c+1)\ge0$. Each summand in Eq.~\ref{eq:u_a} is non negative. Hence $u_a\ge0$.

\subsection{2:4 Feasibility Preservation}
\label{app:24_feasibility}
For each eligible row $i$ and complete valid block $k$, Eq.~\ref{eq:mask} selects the two local indices in $I_{i,k}^{(l)}$ and assigns mask value $1$ only to those entries. The remaining two entries receive mask value $0$. Auxiliary padding entries may be used for bookkeeping. They have score $-\infty$ and are not trainable. They cannot be selected and cannot consume active slots. Hence every complete valid 4-block contains exactly two active trainable entries. Coordinates outside complete valid blocks are excluded from the exact 2:4 claim. They are reported through the valid block fraction. They are not folded into the numerator or denominator of the structured sparsity guarantee. The score projection objective decomposes over independent 4-blocks. Selecting the two largest scores in each block solves the projection exactly.

\subsection{Masked Gradient Isolation}
For each entry, $W_{\text{eff},ij}=W_{ij}M_{ij}$ with the mask fixed during the mini batch. By the chain rule,
\begin{equation}
    \frac{\partial L}{\partial W_{ij}}
    =
    \frac{\partial L}{\partial W_{\text{eff},ij}}
    \frac{\partial W_{\text{eff},ij}}{\partial W_{ij}}
    =
    \frac{\partial L}{\partial W_{\text{eff},ij}}M_{ij}.
\end{equation}
Stacking this entry wise identity over the tensor gives Eq.~\ref{eq:masked_gradient}. If $M_{ij}=0$, the instantaneous task gradient with respect to $W_{ij}$ is zero, which motivates the separate structural proxy pass used by the pruning/regrowth update.

\subsection{Virtual Effective Weight Probe for Dormant Saliency}
\label{app:virtual_probe_note}
The dormant link saliency used by regrowth is not the ordinary parameter gradient $\partial L/\partial W_{ij}$ under a fixed mask. Instead, the structural refresh introduces an auxiliary dense effective tensor $\widetilde{\bm{W}}_{\text{eff}}$ initialized at $\bm{W}\odot\bm{M}$. It computes proxy gradients with respect to that tensor. For a dormant coordinate, $\widetilde{W}_{\text{eff},ij}=0$. The quantity $\partial L_{\text{grow}}/\partial \widetilde{W}_{\text{eff},ij}$ can still be nonzero. The coordinate is treated as a virtual effective input to the proxy computation. These gradients are discarded after ranking. They are never applied as ordinary updates to dormant parameters. Masked gradient isolation for training and virtual probe saliency for topology scoring are compatible.

\subsection{Signed First Order Pruning Criterion}
\label{app:signed_pruning}
Let the mask be fixed during the structural proxy pass and consider an active effective coordinate $w_{\text{eff},ij}$. Define the removal perturbation $\delta_{ij}=-w_{\text{eff},ij}$. Taylor's theorem gives
\begin{equation}
    R_{\mathcal B}(w_{\text{eff},ij}+\delta_{ij})-R_{\mathcal B}(w_{\text{eff},ij})
    =
    \delta_{ij}\frac{\partial R_{\mathcal B}}{\partial w_{\text{eff},ij}}
    +
    \mathcal{R}_{ij}.
\end{equation}
If the one dimensional second derivative along this coordinate is bounded by $L_{ij}$ on the segment between $w_{\text{eff},ij}$ and $0$, then
\[
    |\mathcal{R}_{ij}|
    \le
    \frac{L_{ij}}{2}|w_{\text{eff},ij}|^2.
\]
This gives
\[
    R_{\mathcal B}(w_{\text{eff},ij}=0)-R_{\mathcal B}(w_{\text{eff},ij})
    =
    -w_{\text{eff},ij}\frac{\partial R_{\mathcal B}}{\partial w_{\text{eff},ij}}
    +
    \mathcal{R}_{ij}.
\]
For active links, $w_{\text{eff},ij}=w_{ij}$. The leading term is negative when
\begin{equation}\label{eq:proof_signed_positive}
    w_{\text{eff},ij}\frac{\partial R_{\mathcal B}}{\partial w_{\text{eff},ij}}>0.
\end{equation}
If the leading term dominates the curvature remainder, removal decreases the proxy risk ratio. Conversely, when
\begin{equation}\label{eq:proof_signed_negative}
    w_{\text{eff},ij}\frac{\partial R_{\mathcal B}}{\partial w_{\text{eff},ij}}<0,
\end{equation}
the leading term is positive and removal increases the risk ratio to first order. Otherwise, the sign remains a local saliency ranking. It is not a deterministic descent certificate. The positive part operator in Eq.~\ref{eq:C_ij} keeps only the first case as a pruning candidate.

\subsection{Local Regrowth Saliency}
\label{app:regrowth_saliency}
Let $Q(\widetilde{\bm{W}}_{\text{eff}})=L_{\text{grow}}(\widetilde{\bm{W}}_{\text{eff}})$ denote the structural probe objective evaluated at the current masked effective tensor. For a scalar virtual perturbation $\delta_{ij}$ at connection $(i,j)$,
\begin{equation}
    Q(\widetilde{w}_{\text{eff},ij}+\delta_{ij})-Q(\widetilde{w}_{\text{eff},ij})
    =
    \delta_{ij}\frac{\partial Q}{\partial \widetilde{w}_{\text{eff},ij}}
    +
    \mathcal{R}^{Q}_{ij}.
\end{equation}
If the coordinate wise curvature is bounded by $H_{ij}$ for $|\delta_{ij}|\le\eta$, then
\[
    |\mathcal{R}^{Q}_{ij}|
    \le
    \frac{H_{ij}}{2}\eta^2.
\]
Thus, to leading order, the largest achievable absolute perturbation under $|\delta_{ij}|\le\eta$ is
\[
    \eta\left|
    \frac{\partial Q}{\partial \widetilde{w}_{\text{eff},ij}}
    \right|.
\]
Ranking dormant candidates by Eq.~\ref{eq:G_ij} selects links with high local leading order influence on the KL uniform structural saliency objective. The sign is discarded. The score is not a descent direction. It is used only for candidate topology exploration before the valid block 2:4 projection.

\subsection{Bounded Morphological State}
The final operation in Eq.~\ref{eq:score_update} applies an element wise clamp to the interval $[-V_{\max},V_{\max}]$. Hence every updated score lies in that interval regardless of the magnitude of the unclamped momentum step. The mask is then recomputed from the bounded score tensor using Eq.~\ref{eq:mask}. Appendix~\ref{app:24_feasibility} applies at the next refresh.

\subsection{Controlled Imbalance Modulation}
Because $\operatorname{sg}[\hat{p}_{y_n}^{(n)}]$ is treated as constant during backpropagation, $\nabla_{\bm{e}}\mathcal{F}_n(t)=0$. Thus the sample level gradient is proportional to
\begin{equation}
    \mathcal{F}_n(t)
    \nabla_{\bm{e}}\!\left(
    \bar{\omega}_{y_n}L_{\text{CE}}^{(n)}
    +\lambda_{\text{KL}}\mu(t)
    L_{\text{KL}}^{(n)}
    \right).
\end{equation}
Moreover, $\mathcal{F}_n(t)\ge0$, $\mu(t)\in[0,1]$, and $\bar{\omega}_{y_n}\le\omega_{\max}$ by Eq.~\ref{eq:sample_weight}. In the reported symmetric KL setting, the KL term has no additional class dependent multiplier. In the asymmetric diagnostic variant, the multiplier is capped at $\Lambda_{\max}$ by Eq.~\ref{eq:asym_multiplier}. Hence focal modulation changes sample emphasis without adding an uncontrolled confidence derivative path. The supervised class weight is explicitly capped.

\FloatBarrier

\section{Novelty Boundary Against Closest Prior Art}
\label{app:novelty}
This appendix positions DST-EDL against adjacent method families. The purpose
is to make the novelty boundary explicit. It does not assert general
superiority over alternatives.

\begin{table}[!t]
\centering
\caption{Compact novelty boundary against closest related methods. The distinction is Dirichlet derived proxy saliency used directly for dynamic valid block 2:4 topology adaptation with dormant link structural scoring.}
\label{tab:novelty}
\scriptsize
\renewcommand{\arraystretch}{1.05}
\setlength{\tabcolsep}{2pt}
\begin{adjustbox}{max width=\columnwidth}
\begin{tabular}{@{}p{0.28\columnwidth}p{0.64\columnwidth}@{}}
\toprule
\textbf{Method family} & \textbf{Boundary vs. DST-EDL} \\
\midrule
SNIP~\cite{lee2019snip} & One shot connection sensitivity at initialization. No uncertainty guided dynamic topology. \\
SET~\cite{mocanu2018scalable} / SNFS~\cite{dettmers2019sparse} & Dynamic sparse training with random or task gradient driven regrowth. No Dirichlet ambiguity/vacuity driven regrowth. \\
RigL~\cite{evci2020rigging} / SRigL~\cite{lasby2023srigl} & Strong dynamic sparse baselines. No evidential ambiguity/vacuity controller for exact valid block 2:4 adaptation. \\
Fisher EDL~\cite{deng2023fisher}, R-EDL/Re-EDL~\cite{chen2023redl,chen2024reedl}, Flexible EDL~\cite{yoon2025flexible} & Refine evidential uncertainty modeling. No pruning, regrowth, or structured sparse feasibility. \\
NVIDIA 2:4 sparsity~\cite{mishra2021accelerating} & Provides a hardware compatible sparsity constraint, not uncertainty aware topology learning. \\
\textbf{DST-EDL (ours)} & \textbf{Uses Dirichlet ambiguity/vacuity and KL uniform proxy gradients for pruning/regrowth under exact valid block 2:4 feasibility, including dormant link structural scoring.} \\
\bottomrule
\end{tabular}
\end{adjustbox}
\end{table}

\FloatBarrier

\section{Additional Diagnostics and Future Validation Protocols}
\label{app:external}
The completed softmax and CIFAR-100-LT diagnostics below are included as
additional scope evidence. Broader validation on additional imbalance
profiles, backbones, sparse Tensor Core compatible kernels, and quality gated
failure detection remains outside the central ISIC claim.

\subsection{ISIC Softmax Long Tailed Diagnostic}
Table~\ref{tab:isic_softmax_lt_diagnostic} reports the completed three seed
softmax long tailed baselines on the same ISIC split. These runs are useful for
context. They are not a replacement for the main evidential comparison. They use
ordinary softmax probabilities. They do not expose Dirichlet vacuity, expected
entropy, or valid block sparse topology signals. The strongest softmax row in
this diagnostic is MiSLAS, with pAUC$_{0.80}$ $0.1032\pm0.0064$ and
Macro AUROC $0.8708\pm0.0090$, below the reported DST-EDL means of
$0.1168$ and $0.8891$ under the evidential sparse protocol.

The diagnostic also shows a thresholding pathology typical of extreme prevalence
screening. Standard CE, Focal Loss, LDAM-DRW, cRT, and MiSLAS produce no
positive predictions at the default threshold across all three seeds
(default positive rate $0.0000$). Their default threshold balanced accuracy
collapses to $0.5000$ despite nontrivial ranking metrics. Class Balanced CE and
Balanced Softmax predict a small positive fraction at the default threshold. They
remain substantially lower in pAUC$_{0.80}$ and Macro AUROC. These results motivate the
paper's use of ranking metrics and validation selected operating points rather
than default threshold accuracy under the ISIC 2024 prevalence. The very low
softmax ECE values should also be interpreted with this prevalence effect in
mind: they primarily measure probability calibration under a negative dominated
test distribution, not evidential uncertainty quality.

\begin{table}[!t]
\centering
\scriptsize
\caption{ISIC 2024 softmax long tailed diagnostic over seeds 42, 123, and 456. Values are mean $\pm$ std. Default+ is the default threshold positive prediction rate.}
\label{tab:isic_softmax_lt_diagnostic}
\setlength{\tabcolsep}{1.5pt}
\begin{adjustbox}{max width=\columnwidth}
\begin{tabular}{@{}l C C C C C@{}}
\toprule
Model & pAUC$_{0.80}$ & Macro AUROC & PR AUC & ECE & Default+ \\
\midrule
Standard CE & $0.0927\pm0.0112$ & $0.8561\pm0.0109$ & $0.0323\pm0.0059$ & $0.0007\pm0.0000$ & $0.0000\pm0.0000$ \\
Focal Loss & $0.0941\pm0.0137$ & $0.8585\pm0.0213$ & $0.0363\pm0.0082$ & $0.0015\pm0.0002$ & $0.0000\pm0.0000$ \\
Class Balanced CE & $0.0729\pm0.0107$ & $0.8314\pm0.0147$ & $0.0301\pm0.0036$ & $0.0028\pm0.0003$ & $0.0047\pm0.0002$ \\
Balanced Softmax & $0.0743\pm0.0119$ & $0.8366\pm0.0101$ & $0.0314\pm0.0034$ & $0.0017\pm0.0002$ & $0.0031\pm0.0003$ \\
LDAM-DRW & $0.0985\pm0.0063$ & $0.8623\pm0.0046$ & $0.0319\pm0.0052$ & $0.0009\pm0.0001$ & $0.0000\pm0.0000$ \\
cRT & $0.0905\pm0.0078$ & $0.8499\pm0.0137$ & $\mathbf{0.0424\pm0.0148}$ & $0.0009\pm0.0001$ & $0.0000\pm0.0000$ \\
MiSLAS & $\mathbf{0.1032\pm0.0064}$ & $\mathbf{0.8708\pm0.0090}$ & $0.0313\pm0.0051$ & $0.0015\pm0.0001$ & $0.0000\pm0.0000$ \\
\bottomrule
\end{tabular}
\end{adjustbox}
\end{table}

\subsection{CIFAR-100-LT Diagnostic and Protocol}
Beyond the completed 1:100 diagnostic reported here, broader long tailed
evaluation should vary imbalance profiles (1:10, 1:50, 1:100) and report
top 1/top 5 accuracy, balanced accuracy, macro F1, many/medium/few shot accuracy,
macro AUROC, macro PR AUC, NLL, Brier score, ECE, AURC, and failure detection
metrics against standard long tailed methods (CE, Focal Loss, Logit Adjustment,
Class Balanced Loss, Balanced Softmax, LDAM-DRW, cRT, and MiSLAS style LAS+cRT),
evidential baselines (Dense EDL, Fisher EDL, Flexible EDL, and R-EDL), and
sparse baselines (Static 2:4 EDL and RigL style dynamic sparsity).

A completed CIFAR-100-LT 1:100 run over seeds 42, 123, and 456 is included only as a diagnostic stress test. The diagnostic table below shows that DST-EDL remains structurally healthy and obtains lower ECE than the dense CIFAR baselines. It does not outperform Standard CE or Dense EDL on classification or ranking metrics. This negative result supports the paper's conservative scope. The present contribution is a well defined uncertainty guided valid block sparse mechanism validated primarily in the ISIC screening setting. It is not a universal long tailed classification recipe.

\begin{table}[!t]
\centering
\scriptsize
\caption{CIFAR-100-LT 1:100 diagnostic over seeds 42, 123, and 456. Values are mean $\pm$ std. The run is reported as a scope and stress test result, not as a central performance claim.}
\label{tab:cifar_ir100_diagnostic}
\setlength{\tabcolsep}{1.8pt}
\begin{adjustbox}{max width=\columnwidth}
\begin{tabular}{@{}l C C C C@{}}
\toprule
Model & Bal. & Macro AUROC & ECE & AURC \\
\midrule
Standard CE & $0.4097\pm0.0061$ & $\mathbf{0.9409\pm0.0029}$ & $0.2578\pm0.0116$ & $\mathbf{0.3513\pm0.0093}$ \\
Dense EDL & $0.3623\pm0.0148$ & $0.7600\pm0.0080$ & $0.2321\pm0.0216$ & $0.4809\pm0.0409$ \\
DST-EDL & $0.3411\pm0.0080$ & $0.7524\pm0.0059$ & $\mathbf{0.2036\pm0.0074}$ & $0.5151\pm0.0170$ \\
\bottomrule
\end{tabular}
\end{adjustbox}
\end{table}

For DST-EDL, all three CIFAR-100-LT runs preserved 50.0\% active density, satisfied all valid 2:4 blocks in the wrapped layers (2,789,376/2,789,376), and showed no detected representational collapse. The long tail breakdown remains few shot limited. Mean many shot, medium shot, and few shot accuracies are 0.6004, 0.3063, and 0.0864, respectively. The diagnostic separates structural feasibility from task level competitiveness. Exact 2:4 validity remains intact. CIFAR-100-LT performance indicates that additional tuning or dataset specific design is needed before making broader long tailed claims.

\subsection{Hardware Profiling Protocol}
Hardware evaluation should separate three claims. The first claim is structural feasibility. The second claim is masked PyTorch execution cost. The third claim is realized sparse kernel acceleration. Prior 2:4 work reports that Sparse Tensor Cores can increase matrix math throughput under strict 50\% structured sparsity~\cite{mishra2021accelerating}. Deployment libraries such as cuSPARSELt target matrix multiplication with a 50\% structured sparse operand~\cite{nvidiacusparselt}. PyTorch semi structured tensors similarly require CUDA tensors with supported shape and dtype constraints~\cite{pytorchsemistructured}. Therefore, valid masks alone are not sufficient evidence of speedup.

\begin{table}[!t]
\centering
\scriptsize
\caption{Recommended hardware evidence ladder. Each level answers a different reviewer question. Real acceleration is claimed only at Level 3.}
\label{tab:hardware_protocol_ladder}
\setlength{\tabcolsep}{1.4pt}
\renewcommand{\arraystretch}{1.08}
\begin{adjustbox}{max width=\columnwidth}
\begin{tabular}{@{}p{0.16\columnwidth}p{0.28\columnwidth}p{0.31\columnwidth}p{0.20\columnwidth}@{}}
\toprule
\textbf{Level} & \textbf{Comparison} & \textbf{Measurements} & \textbf{Allowed claim} \\
\midrule
Structural & Dense EDL, Static 2:4, RigL style 2:4, DST-EDL & Active density, valid 2:4 blocks, wrapped layer parameter reduction, theoretical MAC and FLOP reduction & Sparse compatibility \\
Masked PyTorch & Dense EDL, Static 2:4, DST-EDL & Images per second, milliseconds per batch, peak CUDA memory, batch size, image size, GPU name & Diagnostic cost only \\
Sparse kernel & Dense exported model vs. exported 2:4 model & cuSPARSELt, TensorRT, or PyTorch semi structured latency; throughput; memory; layout conversion status & Real speedup only if backend uses sparse kernels \\
\bottomrule
\end{tabular}
\end{adjustbox}
\end{table}

The current implementation can report Levels 1 and 2 with \texttt{experiments/hardware\_profile.py}. The script records active density, valid 2:4 block fraction, dense execution MACs, theoretical active MACs, FLOPs under the two MACs per multiply add convention, masked PyTorch throughput, latency, and peak CUDA memory. These numbers should be presented as structural evidence and masked PyTorch diagnostics. If masked sparse inference is slower than dense inference, the result is still informative. Masked kernels multiply by dense tensors and do not expose the compressed 2:4 layout to Sparse Tensor Cores. Level 3 requires exporting frozen valid 2:4 weights to a compatible backend. If convolutional layers cannot be mapped to supported sparse kernels, the paper should report the failure mode rather than infer acceleration from the mask.

\subsection{Quality Gated Failure Detection Protocol}
A quality control filter based on the entropy based ambiguity proxy $u_a$ is planned for future evaluation. High $u_a$ scores will be tested for correlation with input artifacts such as camera glare, motion blur, and obscurations. Evaluating accepted vs. deferred subsets based on this uncertainty threshold will demonstrate selective prediction utility.

For controlled long tailed recognition, CIFAR-100-LT should be used to separate general imbalance behavior from the ISIC screening setting. Hardware profiling should distinguish structural 2:4 feasibility from realized acceleration. Actual speedup requires sparse Tensor Core compatible kernels and layouts such as cuSPARSELt or TensorRT sparse execution. Quality gated failure detection based on $u_a$ or related ambiguity scores is left as future work. It should be evaluated as selective prediction. It should not be evaluated as a clinical deployment claim.

\end{document}
